\hojaejercicio{Ecuaciones Lineales de Primer Orden y Variación de Constantes}

% Posibilidades para los ejercicios:
% enunciado, medio, resuelto
% profesor, no

% Ejercicio 1
\ejercicio{resuelto}{profesor}{
    Demuestra el siguiente principio de superposición generalizado. Sea $m \geq 2$ un número entero, y
    \[
    A(t) := (a_{ij}(t))_{1 \leq i,j \leq N} \quad \text{y} \quad B_k(t) := (b_{i,k}(t))_{1 \leq i \leq N, \atop 1 \leq k \leq m}
    \]
    Si $u_k$, $1 \leq k \leq m$, son, respectivamente, las soluciones de los sistemas lineales
    \[
    u' = A(t)u + B_k(t),
    \]
    entonces
    \[
    u(t) := \sum_{k=1}^{m} u_k(t)
    \]
    resuelve el sistema lineal
    \[
    u' = A(t)u + \sum_{k=1}^{m} B_k(t).
    \]
}{
    Para demostrar el principio de superposición generalizado, definimos la función $u(t)$ como la suma de las soluciones individuales $u_k(t)$ para $k = 1, \dots, m$:
    \[
    u(t) := \sum_{k=1}^{m} u_k(t).
    \]
    Nuestro objetivo es verificar que esta función $u(t)$ satisface la ecuación diferencial lineal cuyo término no homogéneo es la suma de los términos $B_k(t)$. Procedemos derivando $u(t)$ respecto a $t$. Utilizando la linealidad del operador derivada, tenemos:
    \[
    u'(t) = \frac{d}{dt} \left( \sum_{k=1}^{m} u_k(t) \right) = \sum_{k=1}^{m} u_k'(t).
    \]
    Por hipótesis, sabemos que cada función $u_k(t)$ es solución del sistema lineal correspondiente, lo que significa que satisface la identidad $u_k'(t) = A(t)u_k(t) + B_k(t)$ para todo $1 \leq k \leq m$. Sustituyendo estas expresiones en la sumatoria anterior, obtenemos:
    \begin{align*}
    u'(t) &= \sum_{k=1}^{m} \left( A(t)u_k(t) + B_k(t) \right) \\
        &= \sum_{k=1}^{m} A(t)u_k(t) + \sum_{k=1}^{m} B_k(t).
    \end{align*}
    En virtud de la propiedad distributiva del producto matricial respecto a la suma vectorial, el operador lineal $A(t)$ puede extraerse del primer sumatorio. Esto nos permite reescribir el término como:
    \[
    \sum_{k=1}^{m} A(t)u_k(t) = A(t) \left( \sum_{k=1}^{m} u_k(t) \right) = A(t)u(t).
    \]
    Finalmente, al reintegrar este resultado en la ecuación desarrollada, concluimos que:
    \[
    u'(t) = A(t)u(t) + \sum_{k=1}^{m} B_k(t).
    \]
    Queda así demostrado que la suma $u(t)$ es solución del sistema lineal con el término forzante agregado. 
    \qed
}

% Ejercicio 2
\ejercicio{resuelto}{profesor}{
    Encuentra la solución general de la ecuación lineal escalar
    \[
    u' - 2tu = t.
    \]
}{
    Consideramos la ecuación sin el término no homogéneo:
    \[
    u_h' - 2tu_h = 0.
    \]
    Separamos variables para resolverla:
    \[
    \frac{u_h'}{u_h} = 2t \implies \int \frac{du_h}{u_h} = \int 2t \, dt \implies \ln|u_h| = t^2.
    \]
    Tomando exponenciales, obtenemos la solución homogénea básica (ignorando por un momento la constante multiplicativa):
    \[
    u_h(t) = e^{t^2}.
    \]
    Buscamos una solución de la forma $u(t) = C(t) \cdot e^{t^2}$, donde $C(t)$ es una función a determinar. Derivamos esta expresión usando la regla del producto:
    \[
    u'(t) = C'(t)e^{t^2} + C(t) \cdot (2te^{t^2}).
    \]
    Sustituimos $u(t)$ y $u'(t)$ en la ecuación diferencial original $u' - 2tu = t$:
    \[
    \left[ C'(t)e^{t^2} + 2tC(t)e^{t^2} \right] - 2t \left[ C(t)e^{t^2} \right] = t.
    \]
    Observamos que los términos que contienen $C(t)$ se cancelan (propiedad fundamental del método):
    \[
    C'(t)e^{t^2} = t.
    \]
    Despejamos $C'(t)$:
    \[
    C'(t) = t e^{-t^2}.
    \]

    Integramos respecto a $t$ (usando el cambio de variable $v = -t^2 \implies dv = -2t dt$):
    \[
    C(t) = \int t e^{-t^2} \, dt = -\frac{1}{2} \int e^v \, dv = -\frac{1}{2}e^{-t^2} + K,
    \]
    donde $K$ es una constante arbitraria real.

    Sustituimos $C(t)$ en la expresión propuesta en el Paso 2:
    \[
    u(t) = \left( -\frac{1}{2}e^{-t^2} + K \right) e^{t^2}.
    \]
    Distribuyendo el término exponencial:
    \[
    u(t) = -\frac{1}{2} + K e^{t^2}.
    \]

}

% Ejercicio 3
\ejercicio{resuelto}{profesor}{
    Encuentra la solución única del problema de valor inicial
    \[
    \begin{cases}
    u' + 2tu = t, & t \in \mathbb{R}, \\
    u(1) = 2.
    \end{cases}
    \]
}{
    Primero, hallamos la solución de la ecuación homogénea asociada, ignorando temporalmente el término forzante:
    \[
    u_h' + 2tu_h = 0.
    \]
    Separamos las variables para resolverla:
    \[
    \frac{u_h'}{u_h} = -2t \implies \int \frac{du_h}{u_h} = \int -2t \, dt \implies \ln|u_h| = -t^2.
    \]
    Tomando exponenciales a ambos lados, obtenemos la solución homogénea básica:
    \[
    u_h(t) = e^{-t^2}.
    \]
    
    Para hallar la solución general, aplicamos el método de variación de las constantes. Buscamos una solución de la forma $u(t) = C(t) \cdot e^{-t^2}$, donde $C(t)$ es una función a determinar. Derivamos esta expresión utilizando la regla del producto:
    \[
    u'(t) = C'(t)e^{-t^2} + C(t) \cdot (-2te^{-t^2}).
    \]
    Sustituimos $u(t)$ y $u'(t)$ en la ecuación diferencial original $u' + 2tu = t$:
    \[
    \left[ C'(t)e^{-t^2} - 2tC(t)e^{-t^2} \right] + 2t \left[ C(t)e^{-t^2} \right] = t.
    \]
    Observamos que los términos que contienen $C(t)$ se cancelan mutuamente, lo cual verifica la consistencia del método:
    \[
    C'(t)e^{-t^2} = t.
    \]
    Despejamos la derivada $C'(t)$ multiplicando por $e^{t^2}$:
    \[
    C'(t) = t e^{t^2}.
    \]
    
    Ahora integramos respecto a $t$ para encontrar $C(t)$. Utilizamos el cambio de variable $w = t^2$, lo que implica $dw = 2t \, dt$, o equivalentemente, $\frac{1}{2}dw = t \, dt$:
    \[
    C(t) = \int t e^{t^2} \, dt = \frac{1}{2} \int e^w \, dw = \frac{1}{2}e^{t^2} + K,
    \]
    donde $K$ es la constante arbitraria de integración.
    
    Sustituimos la función $C(t)$ encontrada en nuestra propuesta inicial:
    \[
    u(t) = \left( \frac{1}{2}e^{t^2} + K \right) e^{-t^2}.
    \]
    Distribuyendo el término exponencial $e^{-t^2}$, obtenemos la solución general explícita:
    \[
    u(t) = \frac{1}{2} + K e^{-t^2}.
    \]
    
    Finalmente, para encontrar la solución única del problema, imponemos la condición inicial $u(1) = 2$:
    \[
    2 = \frac{1}{2} + K e^{-(1)^2} \implies 2 - \frac{1}{2} = K e^{-1} \implies \frac{3}{2} = \frac{K}{e}.
    \]
    Despejamos $K$:
    \[
    K = \frac{3e}{2}.
    \]
    Sustituyendo este valor de $K$ en la solución general, concluimos con la solución del problema de valor inicial:
    \[
    u(t) = \frac{1}{2} + \left( \frac{3e}{2} \right) e^{-t^2} = \frac{1}{2} + \frac{3}{2}e^{1-t^2}.
    \]
}

% Ejercicio 4
\ejercicio{resuelto}{profesor}{
    Encuentra la solución única del problema de valor inicial
    \[
    \begin{cases}
    u' = \frac{u}{t} + e^t, & t \in (0, \infty), \\
    u(1) = 2.
    \end{cases}
    \]
}{
    Reescribimos la ecuación diferencial en su forma estándar $u' - \frac{1}{t}u = e^t$. Para aplicar el método de variación de las constantes, comenzamos resolviendo la ecuación homogénea asociada $u_h' - \frac{1}{t}u_h = 0$. Separando variables e integrando:
    \[
    \int \frac{du_h}{u_h} = \int \frac{dt}{t} \implies \ln|u_h| = \ln|t|.
    \]
    Dado que el dominio es $t > 0$, podemos prescindir del valor absoluto, obteniendo la solución homogénea fundamental $u_h(t) = t$.
    
    Proponemos entonces una solución general de la forma $u(t) = C(t)t$, donde $C(t)$ es la función incógnita. Calculamos su derivada mediante la regla del producto:
    \[
    u'(t) = C'(t)t + C(t).
    \]
    Sustituimos $u(t)$ y $u'(t)$ en la ecuación diferencial original $u' = \frac{u}{t} + e^t$:
    \[
    C'(t)t + C(t) = \frac{C(t)t}{t} + e^t.
    \]
    Simplificando el lado derecho, tenemos $C'(t)t + C(t) = C(t) + e^t$. Cancelamos el término $C(t)$ en ambos lados, lo que confirma la validez del planteamiento:
    \[
    C'(t)t = e^t \implies C'(t) = \frac{e^t}{t}.
    \]
    Para hallar $C(t)$, integramos la expresión anterior. Dado que la integral $\int \frac{e^t}{t} dt$ no posee una primitiva expresable mediante funciones elementales, utilizamos el Teorema Fundamental del Cálculo para expresar la solución en términos de una integral definida, tomando $t_0 = 1$ (el punto de la condición inicial) como límite inferior para facilitar el cálculo posterior:
    \[
    C(t) = \int_{1}^{t} \frac{e^s}{s} \, ds + K,
    \]
    donde $s$ es una variable muda de integración y $K$ es una constante arbitraria. La solución general es entonces:
    \[
    u(t) = t \left( K + \int_{1}^{t} \frac{e^s}{s} \, ds \right).
    \]
    Finalmente, aplicamos la condición inicial $u(1) = 2$. Evaluamos la solución en $t=1$:
    \[
    2 = 1 \cdot \left( K + \int_{1}^{1} \frac{e^s}{s} \, ds \right).
    \]
    Como la integral definida desde 1 hasta 1 es cero, obtenemos directamente que $K = 2$. Sustituyendo este valor, llegamos a la solución única del problema:
    \[
    u(t) = 2t + t \int_{1}^{t} \frac{e^s}{s} \, ds.
    \]  
}

% Ejercicio 5
\ejercicio{resuelto}{profesor}{
    Encuentra la solución única del problema de valor inicial
    \[
    \begin{cases}
    u' + u = \frac{1}{1+t^2}, & t \in \mathbb{R}, \\
    u(2) = 3.
    \end{cases}
    \]
}{
    Iniciamos el proceso resolviendo la ecuación homogénea asociada $u_h' + u_h = 0$. Mediante separación de variables, tenemos $\frac{u_h'}{u_h} = -1$, cuya integración conduce a $\ln|u_h| = -t$, resultando en la función base $u_h(t) = e^{-t}$.

    Aplicando el método de variación de las constantes, buscamos una solución de la estructura $u(t) = C(t)e^{-t}$. Derivamos esta expresión respecto a $t$:
    \[
    u'(t) = C'(t)e^{-t} - C(t)e^{-t}.
    \]
    Introducimos $u(t)$ y $u'(t)$ en la ecuación no homogénea original:
    \[
    \left[ C'(t)e^{-t} - C(t)e^{-t} \right] + C(t)e^{-t} = \frac{1}{1+t^2}.
    \]
    Tras la cancelación de los términos semejantes, obtenemos la ecuación diferencial para $C(t)$:
    \[
    C'(t)e^{-t} = \frac{1}{1+t^2} \implies C'(t) = \frac{e^t}{1+t^2}.
    \]
    Integramos para encontrar $C(t)$. Al igual que en el ejercicio anterior, la función $\frac{e^t}{1+t^2}$ no tiene una antiderivada elemental. Por conveniencia, definimos la integral acumulada desde el punto inicial $t_0 = 2$:
    \[
    C(t) = \int_{2}^{t} \frac{e^s}{1+s^2} \, ds + K.
    \]
    Sustituyendo $C(t)$ en nuestra propuesta inicial, la solución general toma la forma:
    \[
    u(t) = e^{-t} \left( K + \int_{2}^{t} \frac{e^s}{1+s^2} \, ds \right).
    \]
    Para determinar la constante $K$, imponemos la condición inicial $u(2) = 3$:
    \[
    3 = e^{-2} \left( K + \int_{2}^{2} \frac{e^s}{1+s^2} \, ds \right).
    \]
    La integral se anula, resultando en $3 = e^{-2}K$, de donde despejamos $K = 3e^2$. Reemplazamos $K$ en la expresión general para obtener la solución final:
    \[
    u(t) = e^{-t} \left( 3e^2 + \int_{2}^{t} \frac{e^s}{1+s^2} \, ds \right) = 3e^{2-t} + e^{-t} \int_{2}^{t} \frac{e^s}{1+s^2} \, ds.
    \]    
}

% Ejercicio 6
\ejercicio{resuelto}{profesor}{
    Encuentra la solución única continua del problema de valor inicial
    \[
    \begin{cases}
    u' + u = b(t), \\
    u(0) = 0
    \end{cases}
    \]
    en el intervalo $[0, \infty)$, donde
    \[
    b(t) = \begin{cases} 2 & \text{si } 0 \leq t \leq 1, \\ 0 & \text{si } t > 1. \end{cases}
    \]    
}{
    Dada la naturaleza definida a trozos de la función fuente $b(t)$, debemos resolver la ecuación diferencial por separado en cada intervalo, asegurando posteriormente la continuidad de la función $u(t)$ en el punto de empalme $t=1$.

    Comenzamos analizando el primer intervalo, $0 \leq t \leq 1$, donde $b(t) = 2$. La ecuación diferencial se reduce a $u' + u = 2$. Primero resolvemos la ecuación homogénea asociada $u_h' + u_h = 0$, cuya solución es $u_h(t) = e^{-t}$. Aplicando el método de variación de las constantes, proponemos $u(t) = C(t)e^{-t}$. Al derivar y sustituir en la ecuación completa, obtenemos:
    \[
    C'(t)e^{-t} = 2 \implies C'(t) = 2e^t.
    \]
    Integramos para hallar $C(t) = 2e^t + K_1$. Por consiguiente, la solución general en este intervalo es $u(t) = (2e^t + K_1)e^{-t} = 2 + K_1e^{-t}$. Imponemos la condición inicial global $u(0) = 0$:
    \[
    0 = 2 + K_1e^0 \implies K_1 = -2.
    \]
    Así, la solución para el primer tramo es $u(t) = 2(1 - e^{-t})$. Para garantizar la continuidad, calculamos el valor de la función al final de este intervalo:
    \[
    u(1) = 2(1 - e^{-1}).
    \]
    
    Procedemos ahora con el segundo intervalo, $t > 1$, donde $b(t) = 0$. La ecuación se convierte en homogénea: $u' + u = 0$. La solución general es inmediata: $u(t) = K_2 e^{-t}$. Para determinar la constante $K_2$, utilizamos la condición de continuidad en $t=1$, es decir, $\lim_{t \to 1^-} u(t) = \lim_{t \to 1^+} u(t)$. Igualamos el valor obtenido del primer tramo con la expresión del segundo evaluada en $t=1$:
    \[
    2(1 - e^{-1}) = K_2 e^{-1}.
    \]
    Despejamos $K_2$ multiplicando por $e$:
    \[
    K_2 = 2e(1 - e^{-1}) = 2(e - 1).
    \]
    Sustituyendo este valor, la solución para $t > 1$ es $u(t) = 2(e - 1)e^{-t}$.
    
    En conclusión, la solución única continua del problema de valor inicial es:
    \[
    u(t) = \begin{cases} 
    2(1 - e^{-t}) & \text{si } 0 \leq t \leq 1, \\
    2(e - 1)e^{-t} & \text{si } t > 1.
    \end{cases}
    \]    
}

% Ejercicio 7
\ejercicio{resuelto}{profesor}{
    Sean $a, m, \omega \in \mathbb{R}$ tales que $a > 0$ y $\omega > 0$. Encuentra la solución general de
    \[
    u' = -au + me^{-\omega t}
    \]
    y demuestra que toda solución tiende a cero cuando $t \uparrow \infty$.    
}{
    Reescribimos la ecuación lineal en su forma estándar: $u' + au = me^{-\omega t}$. Resolvemos primero la ecuación homogénea asociada $u_h' + au_h = 0$. Al separar variables e integrar, obtenemos $\ln|u_h| = -at$, lo que nos da la solución fundamental $u_h(t) = e^{-at}$.

    Utilizando el método de variación de las constantes, buscamos una solución de la forma $u(t) = C(t)e^{-at}$. Derivamos esta expresión, $u'(t) = C'(t)e^{-at} - aC(t)e^{-at}$, y la sustituimos en la ecuación original. Tras la cancelación de términos, llegamos a:
    \[
    C'(t)e^{-at} = me^{-\omega t} \implies C'(t) = m e^{(a-\omega)t}.
    \]
    Para integrar esta expresión y hallar $C(t)$, debemos distinguir dos casos según si $a = \omega$ o $a \neq \omega$. Asumiremos el caso general $a \neq \omega$ (el caso $a=\omega$ se resuelve análogamente resultando en un término lineal $t$). Integramos respecto a $t$:
    \[
    C(t) = \int m e^{(a-\omega)t} \, dt = \frac{m}{a-\omega}e^{(a-\omega)t} + K,
    \]
    donde $K$ es una constante arbitraria. Sustituyendo $C(t)$ en la propuesta inicial, la solución general es:
    \[
    u(t) = \left( \frac{m}{a-\omega}e^{(a-\omega)t} + K \right) e^{-at} = \frac{m}{a-\omega}e^{-\omega t} + Ke^{-at}.
    \]
    (Nota: Si $a=\omega$, la solución sería $u(t) = mte^{-at} + Ke^{-at}$).
    
    Finalmente, analizamos el comportamiento asintótico de la solución cuando $t \to \infty$. Observamos la expresión obtenida para el caso general:
    \[
    \lim_{t \to \infty} u(t) = \lim_{t \to \infty} \left( \frac{m}{a-\omega}e^{-\omega t} + Ke^{-at} \right).
    \]
    Dado que por hipótesis los parámetros $a$ y $\omega$ son estrictamente positivos ($a > 0, \omega > 0$), las exponenciales $e^{-\omega t}$ y $e^{-at}$ son decrecientes y tienden a cero cuando $t$ crece indefinidamente. Por lo tanto:
    \[
    \lim_{t \to \infty} u(t) = 0 + 0 = 0.
    \]
    Esta conclusión se mantiene igualmente válida para el caso de resonancia ($a=\omega$), ya que el término $te^{-at}$ también tiende a cero por la regla de L'Hôpital. Queda así demostrado que todas las soluciones tienden a cero asintóticamente.
}

% Ejercicio 8
\ejercicio{resuelto}{profesor}{
    Sean $a, b \in \mathcal{C}(\mathbb{R}; \mathbb{R})$ tales que
    \[
    \sup_{\mathbb{R}} a \leq -c < 0 \quad \text{y} \quad \lim_{t \uparrow \infty} b(t) = 0,
    \]
    para una constante $c$. Demuestra que toda solución de $u' = a(t)u + b(t)$ tiende a cero cuando $t \uparrow \infty$.
}{
    Sea $u(t)$ una solución arbitraria de la ecuación diferencial lineal $u' = a(t)u + b(t)$. Para estudiar su comportamiento asintótico, primero derivamos una expresión explícita para $u(t)$ utilizando el método de variación de las constantes.

    Consideramos un tiempo inicial $t_0$ arbitrario. La solución general de la ecuación lineal de primer orden viene dada por la fórmula:
    \[
    u(t) = u(t_0) \exp\left( \int_{t_0}^t a(s) \, ds \right) + \int_{t_0}^t \exp\left( \int_{\tau}^t a(s) \, ds \right) b(\tau) \, d\tau.
    \]
    Nuestro objetivo es acotar esta expresión utilizando las hipótesis dadas. Sabemos que $\sup_{\mathbb{R}} a \leq -c$, lo que implica que $a(s) \leq -c$ para todo $s \in \mathbb{R}$. En consecuencia, para cualquier intervalo $[\tau, t]$ con $\tau \leq t$, tenemos:
    \[
    \int_{\tau}^t a(s) \, ds \leq \int_{\tau}^t -c \, ds = -c(t - \tau).
    \]
    Utilizando esta desigualdad en la fórmula de la solución y tomando valores absolutos, obtenemos la siguiente acotación para $|u(t)|$:
    \[
    |u(t)| \leq |u(t_0)| e^{-c(t - t_0)} + \int_{t_0}^t e^{-c(t - \tau)} |b(\tau)| \, d\tau.
    \]
    Definimos $I_1(t) = |u(t_0)| e^{-c(t - t_0)}$ e $I_2(t) = \int_{t_0}^t e^{-c(t - \tau)} |b(\tau)| \, d\tau$. Analizaremos el límite de cada término por separado cuando $t \to \infty$.

    \begin{enumerate}
        \item \textbf{Análisis del primer término $I_1(t)$:}
        Dado que $c > 0$, la función exponencial $e^{-ct}$ decrece a cero. Por tanto, independientemente del valor inicial $u(t_0)$, tenemos:
        \[
        \lim_{t \to \infty} I_1(t) = |u(t_0)| e^{ct_0} \lim_{t \to \infty} e^{-ct} = 0.
        \]
        \item \textbf{Análisis del segundo término $I_2(t)$:}
        Debemos demostrar que la integral de convolución tiende a cero. Dado que $\lim_{t \to \infty} b(t) = 0$, para cualquier $\epsilon > 0$, existe un tiempo $T_\epsilon > t_0$ tal que $|b(\tau)| < \frac{c\epsilon}{2}$ para todo $\tau \geq T_\epsilon$.
        Dividimos la integral $I_2(t)$ en dos partes para $t > T_\epsilon$:
        \[
        I_2(t) = \int_{t_0}^{T_\epsilon} e^{-c(t - \tau)} |b(\tau)| \, d\tau + \int_{T_\epsilon}^t e^{-c(t - \tau)} |b(\tau)| \, d\tau.
        \]
        Llamemos $A(t)$ a la primera integral y $B(t)$ a la segunda.
    
        Para $B(t)$, usamos la cota de $|b(\tau)|$:
        \[
        B(t) \leq \int_{T_\epsilon}^t e^{-c(t - \tau)} \frac{c\epsilon}{2} \, d\tau = \frac{c\epsilon}{2} e^{-ct} \int_{T_\epsilon}^t e^{c\tau} \, d\tau = \frac{c\epsilon}{2} e^{-ct} \left[ \frac{e^{c\tau}}{c} \right]_{T_\epsilon}^t.
        \]
        Simplificando:
        \[
        B(t) \leq \frac{\epsilon}{2} e^{-ct} (e^{ct} - e^{cT_\epsilon}) = \frac{\epsilon}{2} (1 - e^{-c(t - T_\epsilon)}) < \frac{\epsilon}{2}.
        \]
        Para $A(t)$, observamos que es una integral sobre un intervalo fijo $[t_0, T_\epsilon]$, por lo que su valor está acotado por una constante $M$. El factor $e^{-c(t-\tau)}$ puede escribirse como $e^{-ct}e^{c\tau}$.
        \[
        A(t) = e^{-ct} \int_{t_0}^{T_\epsilon} e^{c\tau} |b(\tau)| \, d\tau = C_{const} \cdot e^{-ct}.
        \]
        Como $c > 0$, $\lim_{t \to \infty} A(t) = 0$. Por tanto, existe un $T_{final}$ tal que para todo $t > T_{final}$, $A(t) < \frac{\epsilon}{2}$.
    
        Sumando ambas partes, para $t$ suficientemente grande, tenemos $I_2(t) < \frac{\epsilon}{2} + \frac{\epsilon}{2} = \epsilon$. Esto demuestra que $\lim_{t \to \infty} I_2(t) = 0$.
    \end{enumerate}

    Finalmente, como $|u(t)| \leq I_1(t) + I_2(t)$ y ambos términos tienden a cero, concluimos que:
    \[
    \lim_{t \to \infty} u(t) = 0.
    \]
    \qed
}

% Ejercicio 9
\ejercicio{resuelto}{profesor}{
    Sea la siguiente ecuación diferencial:
    \[
        N'(t) = - \lambda N(t)
    \]
    donde $\lambda > 0$ es una constante. Encuentra en que momento $N(t)$ es igual a la mitad de su valor inicial $N(0)$.\\
    Notese que esta ecuación modela el proceso de desintegración radiactiva, y el tiempo que tarda una cantidad de material radiactivo en reducirse a la mitad de su valor inicial se conoce como su \textit{vida media}.
}{
    Se puede ver rapidamente que la solucion general de la ecuacion diferencial es $N(t) = N(0)e^{-\lambda t}$. Para encontrar el tiempo $t_{1/2}$ en el que $N(t)$ es igual a la mitad de su valor inicial, planteamos la siguiente ecuacion:
    \[
        N(0)e^{-\lambda t_{1/2}} = \frac{N(0)}{2} \implies
        e^{-\lambda t_{1/2}} = \frac{1}{2} \implies
        -\lambda t_{1/2} = \ln\left(\frac{1}{2}\right) \implies
        t_{1/2} = \frac{\ln(2)}{\lambda}.
    \]
}

% Ejercicio 10
\ejercicio{resuelto}{profesor}{
    Sean $\alpha < \beta$, $t_0 \in [\alpha, \beta]$, $\eta > 0$ y sea una norma $\|\cdot\|$ en $\mathcal{C}([\alpha, \beta]; \mathbb{R}^N)$. Demuestra que:
    \[
    \|u\|_\eta := \sup_{t \in [\alpha, \beta]} e^{-\eta(t - t_0)} \|u(t)\|
    \]
    define una norma en $\mathcal{C}([\alpha, \beta]; \mathbb{R}^N)$.
}{
    Para demostrar que $\|\cdot\|_\eta$ define una norma, debemos verificar los tres axiomas de norma:
    \begin{enumerate}
        \item \textbf{Positividad:} $\|u\|_\eta \geq 0$ para todo $u \in \mathcal{C}([\alpha, \beta]; \mathbb{R}^N)$ y $\|u\|_\eta = 0$ si y solo si $u = 0$.\\
        Dado que $\|u(t)\| \geq 0$ para todo $t$, y $e^{-\eta(t - t_0)} > 0$, se tiene que $\|u\|_\eta \geq 0$.\\
        Si $\|u\|_\eta = 0$, entonces para todo $t$, $e^{-\eta(t - t_0)} \|u(t)\| = 0$. Como la exponencial es positiva, necesariamente $\|u(t)\| = 0$ para todo $t$, lo cual implica que $u(t) = 0$ para todo $t$.\\
        Recíprocamente, si $u = 0$, entonces $\|u(t)\| = 0$ para todo $t$, por lo tanto $\|u\|_\eta = 0$.
        \item \textbf{Homogeneidad:} Para cualquier escalar $\lambda \in \mathbb{R}$,
        \[
            \|\lambda u\|_\eta = \sup_{t \in [\alpha, \beta]} e^{-\eta(t - t_0)} \|\lambda u(t)\|
            = |\lambda| \sup_{t \in [\alpha, \beta]} e^{-\eta(t - t_0)} \|u(t)\|
            = |\lambda| \|u\|_\eta.
        \]
        Esto se sigue de la propiedad de homogeneidad de la norma en $\mathbb{R}^N$.
        \item \textbf{Desigualdad triangular:} Para cualquier par de funciones $u, v$, tenemos:
        $$
            \|u + v\|_\eta = 
            \sup_{t \in [\alpha, \beta]} e^{-\eta(t - t_0)} \|u(t) + v(t)\|
            \\
            < 
            \sup_{t \in [\alpha, \beta]} e^{-\eta(t - t_0)} (\|u(t)\| + \|v(t)\|)
            \\
            =
            \|u\|_\eta + \|v\|_\eta.
        $$
    \end{enumerate}
    En resumen, todas las propiedades de una norma se cumplen. Por tanto, $\|\cdot\|_\eta$ define una norma en el espacio de funciones continuas sobre el intervalo $[\alpha, \beta]$ con valores en $\mathbb{R}^N$.
}


% Ejercicio 11
\ejercicio{resuelto}{profesor}{
    Para cada entero $n \geq 1$, considere la función $u_n(t) = t^n$, $t \in [0, 1]$. Demuestre que $\{u_n\}_{n \geq 1}$ converge puntualmente en $[0, 1]$ pero no puede converger uniformemente en $[0, 1]$ cuando $n \to \infty$.
}{
    Sea la sucesión de funciones $\{u_n\}_{n \geq 1}$ definida por $u_n(t) = t^n$ en el intervalo $t \in [0, 1]$. Vamos a dividir la demostración en dos partes: primero probaremos la convergencia puntual y luego demostraremos que dicha convergencia no puede ser uniforme.

    Para estudiar la convergencia puntual, debemos fijar un punto $t \in [0, 1]$ y calcular el límite de la sucesión numérica $\{u_n(t)\}$ cuando $n \to \infty$. Distinguimos dos casos según el valor de $t$:

    \begin{enumerate}
        \item \textbf{Caso $0 \leq t < 1$:} En este intervalo, $t$ es un número estrictamente menor que 1. Sabemos que el límite de una progresión geométrica de razón menor que 1 es cero. Por lo tanto:
        $$ \lim_{n \to \infty} u_n(t) = \lim_{n \to \infty} t^n = 0 $$
        
        \item \textbf{Caso $t = 1$:} Si evaluamos la función exactamente en este punto, obtenemos $u_n(1) = 1^n = 1$ para cualquier valor de $n$. En consecuencia, su límite es trivial:
        $$ \lim_{n \to \infty} u_n(1) = \lim_{n \to \infty} 1 = 1 $$
    \end{enumerate}

    Agrupando ambos resultados, concluimos que la sucesión converge puntualmente en todo el intervalo $[0, 1]$ hacia una función límite $u(t)$, definida a trozos de la siguiente manera:
    $$ u(t) = 
    \begin{cases} 
    0 & \text{si } 0 \leq t < 1 \\
    1 & \text{si } t = 1 
    \end{cases} $$

    Para demostrar que esta convergencia no es uniforme, podemos recurrir a dos argumentos distintos y complementarios. Explicaremos ambos para mayor rigor.

    \begin{itemize}
        \item Existe un teorema fundamental en el análisis matemático que establece que si una sucesión de funciones continuas converge uniformemente hacia una función límite, entonces dicha función límite también debe ser continua. 

        En nuestro caso, cada función $u_n(t) = t^n$ es un polinomio y, por tanto, es continua en todo el intervalo $[0, 1]$. Sin embargo, la función límite $u(t)$ que hemos encontrado en el apartado anterior presenta una discontinuidad de salto en $t = 1$. Dado que el límite de una sucesión uniformemente convergente de funciones continuas no puede ser discontinuo, concluimos por reducción al absurdo que la convergencia no puede ser uniforme.

        \item La convergencia uniforme exige que la máxima distancia entre las funciones de la sucesión $u_n(t)$ y la función límite $u(t)$ tienda a cero cuando $n \to \infty$. Matemáticamente, esto significa que el límite del supremo debe ser nulo. Evaluemos dicho supremo:

        $$ M_n = \sup_{t \in [0, 1]} |u_n(t) - u(t)| $$

        Dado que en $t=1$ la diferencia es exacta ($1 - 1 = 0$), el supremo se va a encontrar en el intervalo semiabierto $[0, 1)$. En este intervalo, sabemos que $u(t) = 0$, por lo que la expresión se simplifica a:

        $$ M_n = \sup_{t \in [0, 1)} |t^n - 0| = \sup_{t \in [0, 1)} t^n $$

        A medida que $t$ se acerca a 1 por la izquierda, el valor de $t^n$ (para cualquier $n$ fijo) se acerca arbitrariamente a 1. Por consiguiente, el supremo es exactamente 1:
        $$ M_n = 1 \quad \text{para todo } n \geq 1 $$

        Al tomar el límite cuando $n \to \infty$, obtenemos:
        $$ \lim_{n \to \infty} M_n = \lim_{n \to \infty} 1 = 1 \neq 0 $$

        Como la máxima separación entre $u_n(t)$ y $u(t)$ no tiende a cero, queda demostrado que la convergencia no es uniforme en $[0, 1]$.
    \end{itemize}    
}

% Ejercicio 12
\ejercicio{resuelto}{profesor}{

    Consideramos que en el espcacio vectoy
    Dados $\mathbb{K} \in \{\mathbb{R}, \mathbb{C}\}$ y $\alpha, \beta \in \mathbb{R}$ con $\alpha < \beta$, denotemos por $C^1([\alpha, \beta]; \mathbb{K}^N)$ el espacio vectorial de todas las funciones $u : [\alpha, \beta] \to \mathbb{K}^N$ tales que $u'(t)$ existe para todo $t \in [\alpha, \beta]$ y $u' \in C([\alpha, \beta]; \mathbb{K}^N)$. Demuestre que
    $$ \|u\|_{C^1([\alpha, \beta]; \mathbb{K}^N)} := \|u\|_\infty + \|u'\|_\infty $$
    dota a $C^1([\alpha, \beta]; \mathbb{K}^N)$ con una estructura de espacio de Banach.
}{
    Sea el espacio vectorial $C^1([\alpha, \beta]; \mathbb{K}^N)$ equipado con la norma:
    $$ \|u\|_{C^1} = \|u\|_\infty + \|u'\|_\infty $$
    donde $\|u\|_\infty = \sup_{t \in [\alpha, \beta]} |u(t)|$. Para probar que es un espacio de Banach, demostraremos que toda sucesión de Cauchy en esta norma converge a un elemento del espacio.

    Sea $\{u_n\}_{n \geq 1}$ una sucesión de Cauchy en $C^1([\alpha, \beta]; \mathbb{K}^N)$. Por definición, para cada $\epsilon > 0$, existe $M \in \mathbb{N}$ tal que para todo $n, m > M$:
    $$ \|u_n - u_m\|_{C^1} = \|u_n - u_m\|_\infty + \|u'_n - u'_m\|_\infty < \epsilon $$

    Esto implica directamente que:
    \begin{itemize}
        \item $\|u_n - u_m\|_\infty < \epsilon$, por lo que $\{u_n\}$ es Cauchy en $C([\alpha, \beta]; \mathbb{K}^N)$.
        \item $\|u'_n - u'_m\|_\infty < \epsilon$, por lo que $\{u'_n\}$ es Cauchy en $C([\alpha, \beta]; \mathbb{K}^N)$.
    \end{itemize}

    Dado que el espacio de funciones continuas $(C([\alpha, \beta]; \mathbb{K}^N), \|\cdot\|_\infty)$ es un espacio de Banach completo, existen funciones continuas $u$ y $v$ tales que:
    $$ u_n \xrightarrow{unif.} u \quad \text{y} \quad u'_n \xrightarrow{unif.} v $$

    Debemos verificar que $u$ es derivable y que $u' = v$. Usando el Teorema Fundamental del Cálculo para cada $u_n$:
    $$ u_n(t) = u_n(\alpha) + \int_{\alpha}^{t} u'_n(s) \, ds $$

    Como $u'_n$ converge uniformemente a $v$ en el intervalo compacto $[\alpha, \beta]$, podemos intercambiar el límite con la integral:
    $$ \lim_{n \to \infty} u_n(t) = \lim_{n \to \infty} u_n(\alpha) + \int_{\alpha}^{t} \lim_{n \to \infty} u'_n(s) \, ds $$
    $$ u(t) = u(\alpha) + \int_{\alpha}^{t} v(s) \, ds $$

    Por el Teorema Fundamental del Cálculo, dado que $v$ es continua, la función $u$ es derivable en $[\alpha, \beta]$ y se cumple que $u'(t) = v(t)$. Por lo tanto, $u \in C^1([\alpha, \beta]; \mathbb{K}^N)$.

    Finalmente, evaluamos la convergencia de la sucesión original en la norma del espacio:
    $$ \lim_{n \to \infty} \|u_n - u\|_{C^1} = \lim_{n \to \infty} (\|u_n - u\|_\infty + \|u'_n - u'\|_\infty) $$
    Como $\|u_n - u\|_\infty \to 0$ y $\|u'_n - v\|_\infty \to 0$, tenemos:
    $$ \lim_{n \to \infty} \|u_n - u\|_{C^1} = 0 $$

}


% Ejercicio 13
\ejercicio{resuelto}{profesor}{
    Demuestre que la norma de operador $\|\cdot\|_{\mathcal{L}(\mathbb{K}^N)}$ inducida por una norma $\|\cdot\|$ de $\mathbb{K}^N$, la cual se define como
    $$ \|A\|_{\mathcal{L}(\mathbb{K}^N)} = \max\{\|Ax\| : \|x\| = 1\} $$
    es una norma en el espacio $\mathcal{L}(\mathbb{K}^N)$ de los operadores lineales de $\mathbb{K}^N$ en $\mathbb{K}^N$.
}{
    Para demostrar que la función $\|\cdot\|_{\mathcal{L}(\mathbb{K}^N)}$ define efectivamente una norma en el espacio de operadores lineales $\mathcal{L}(\mathbb{K}^N)$, debemos verificar que cumple con los tres axiomas fundamentales de una norma. Sean $A, B \in \mathcal{L}(\mathbb{K}^N)$ dos operadores lineales cualesquiera y sea $\lambda \in \mathbb{K}$ un escalar.

    \begin{itemize}
        \item \textbf{Positividad:}
        Debemos probar que $\|A\|_{\mathcal{L}(\mathbb{K}^N)} \geq 0$ para todo operador $A$, y que $\|A\|_{\mathcal{L}(\mathbb{K}^N)} = 0$ si y solo si $A = 0$ (el operador nulo).

        En primer lugar, dado que $\|\cdot\|$ es una norma en $\mathbb{K}^N$, sabemos que $\|Ax\| \geq 0$ para cualquier vector $x$. Por consiguiente, el máximo de un conjunto de valores no negativos también es no negativo, lo que garantiza que $\|A\|_{\mathcal{L}(\mathbb{K}^N)} \geq 0$.

        En segundo lugar, supongamos que $\|A\|_{\mathcal{L}(\mathbb{K}^N)} = 0$. Esto significa que:
        $$ \max\{\|Ax\| : \|x\| = 1\} = 0 $$
        Dado que el máximo es cero y todos los valores de $\|Ax\|$ son no negativos, se deduce que $\|Ax\| = 0$ para todo vector $x$ con norma $\|x\| = 1$. Por las propiedades de la norma en $\mathbb{K}^N$, esto implica que $Ax = 0$ para todo vector unitario. 

        Para cualquier vector arbitrario $y \neq 0$ en $\mathbb{K}^N$, podemos construir el vector unitario $u = \frac{y}{\|y\|}$. Aplicando la linealidad del operador $A$:
        $$ A(u) = A\left(\frac{y}{\|y\|}\right) = \frac{1}{\|y\|} Ay = 0 \implies Ay = 0 $$
        Como $A0 = 0$ trivialmente, concluimos que $Ay = 0$ para todo $y \in \mathbb{K}^N$, lo que significa que $A$ es el operador nulo ($A = 0$). 

        La implicación inversa es directa: si $A = 0$, entonces $Ax = 0$ para todo $x$, lo que hace que $\|Ax\| = 0$ y, por tanto, su máximo sobre los vectores unitarios es $0$.

        \item \textbf{Homogeneidad absoluta:}
        Debemos demostrar que $\|\lambda A\|_{\mathcal{L}(\mathbb{K}^N)} = |\lambda| \|A\|_{\mathcal{L}(\mathbb{K}^N)}$.

        Utilizando la definición de la norma de operador y la linealidad de la multiplicación por un escalar, evaluamos:
        $$ \|\lambda A\|_{\mathcal{L}(\mathbb{K}^N)} = \max_{\|x\| = 1} \|(\lambda A)x\| = \max_{\|x\| = 1} \|\lambda (Ax)\| $$
        Dado que $\|\cdot\|$ es una norma en $\mathbb{K}^N$, cumple la propiedad de homogeneidad absoluta, permitiéndonos sacar el escalar en valor absoluto:
        $$ \max_{\|x\| = 1} \|\lambda (Ax)\| = \max_{\|x\| = 1} \left( |\lambda| \|Ax\| \right) $$
        Como $|\lambda|$ es una constante independiente de $x$, podemos extraerla del máximo:
        $$ |\lambda| \max_{\|x\| = 1} \|Ax\| = |\lambda| \|A\|_{\mathcal{L}(\mathbb{K}^N)} $$

        \item \textbf{Desigualdad triangular:}
        Finalmente, debemos probar que $\|A + B\|_{\mathcal{L}(\mathbb{K}^N)} \leq \|A\|_{\mathcal{L}(\mathbb{K}^N)} + \|B\|_{\mathcal{L}(\mathbb{K}^N)}$.

        Por definición, evaluamos la suma de los operadores:
        $$ \|A + B\|_{\mathcal{L}(\mathbb{K}^N)} = \max_{\|x\| = 1} \|(A + B)x\| = \max_{\|x\| = 1} \|Ax + Bx\| $$
        Aplicamos ahora la desigualdad triangular de la norma vectorial subyacente en $\mathbb{K}^N$, la cual establece que $\|Ax + Bx\| \leq \|Ax\| + \|Bx\|$ para cualquier $x$:
        $$ \max_{\|x\| = 1} \|Ax + Bx\| \leq \max_{\|x\| = 1} \left( \|Ax\| + \|Bx\| \right) $$
        Es una propiedad conocida del máximo de funciones que el máximo de una suma es menor o igual a la suma de los máximos. Por lo tanto:
        $$ \max_{\|x\| = 1} \left( \|Ax\| + \|Bx\| \right) \leq \max_{\|x\| = 1} \|Ax\| + \max_{\|x\| = 1} \|Bx\| $$
        Sustituyendo esto por las definiciones de las normas de los operadores individuales, llegamos a:
        $$ \|A + B\|_{\mathcal{L}(\mathbb{K}^N)} \leq \|A\|_{\mathcal{L}(\mathbb{K}^N)} + \|B\|_{\mathcal{L}(\mathbb{K}^N)} $$
    \end{itemize}    
}

% Ejercicio 14
\ejercicio{resuelto}{profesor}{
    Demuestre que, si $\|\cdot\|_{\mathcal{L}(\mathbb{K}^N)}$ es la norma de operador inducida por la norma $p$ de $\mathbb{K}^N$, en general, se tiene que $\|A\|_{\mathcal{L}(\mathbb{K}^N)} \neq \|A\|_p$, donde en esta última norma, se considera $A$ como un vector de $\mathbb{K}^{N^2}$.
}{
    Para demostrar que una afirmación matemática no se cumple "en general", es suficiente con exhibir un contraejemplo donde la igualdad falle. 

    El contraejemplo más natural y elegante para ilustrar la diferencia entre estas dos normas es utilizar la \textbf{matriz identidad} $I_N$ de tamaño $N \times N$, asumiendo un espacio de dimensión $N \geq 2$. Vamos a evaluar ambas normas para esta matriz en el caso general donde el parámetro $p$ cumple $1 \leq p < \infty$.

    Por definición, la norma de operador inducida por una norma $p$ vectorial se define como el máximo estiramiento que la matriz le provoca a un vector de norma 1:
    $$ \|I_N\|_{\mathcal{L}(\mathbb{K}^N)} = \max_{\substack{x \in \mathbb{K}^N \\ \|x\|_p = 1}} \|I_N x\|_p $$

    Dado que la matriz identidad es el operador neutral y no altera a ningún vector ($I_N x = x$), la expresión se simplifica inmediatamente:
    $$ \|I_N\|_{\mathcal{L}(\mathbb{K}^N)} = \max_{\|x\|_p = 1} \|x\|_p = 1 $$

    Por lo tanto, la norma de operador de la matriz identidad es siempre exactamente 1, independientemente de la dimensión $N$ o del valor de $p$.

    Ahora cambiaremos de perspectiva y trataremos a la matriz $I_N$ no como un operador, sino como un único vector "plano" que pertenece al espacio $\mathbb{K}^{N^2}$. 

    Si observamos las componentes de la matriz identidad, esta posee exactamente $N$ elementos en su diagonal principal que tienen un valor de $1$, mientras que los restantes $N^2 - N$ elementos fuera de la diagonal valen $0$. 

    La norma $p$ de este vector de $N^2$ componentes se calcula sumando la potencia $p$-ésima del valor absoluto de todas sus entradas y extrayendo la raíz $p$-ésima de dicho total:
    $$ \|I_N\|_p = \left( \sum_{i=1}^N \sum_{j=1}^N |(I_N)_{i,j}|^p \right)^{\frac{1}{p}} $$

    Desglosando los ceros y los unos, obtenemos:
    $$ \|I_N\|_p = \left( \underbrace{1^p + 1^p + \dots + 1^p}_{N \text{ elementos}} + \underbrace{0^p + \dots + 0^p}_{N^2-N \text{ elementos}} \right)^{\frac{1}{p}} $$
    $$ \|I_N\|_p = (N)^{\frac{1}{p}} = N^{\frac{1}{p}} $$

    Finalmente, comparamos los dos resultados obtenidos para cualquier dimensión $N \geq 2$ y para cualquier norma con $1 \leq p < \infty$. Resulta evidente que:
    $$ N^{\frac{1}{p}} > 1 $$

    Sustituyendo por las normas que hemos calculado:
    $$ \|I_N\|_{\mathcal{L}(\mathbb{K}^N)} = 1 \neq N^{\frac{1}{p}} = \|I_N\|_p $$

    Dado que hemos encontrado una matriz (la identidad) para la cual los valores difieren, queda demostrado que, en general, la norma de operador no es igual a la norma de la matriz tratada como un vector.

    \vspace{0.3cm}
    \textbf{Nota sobre el caso $p = \infty$:} El contraejemplo de la matriz identidad no sirve si $p = \infty$, ya que en ese caso límite $N^{1/\infty} = N^0 = 1$, y las normas coincidirían por casualidad. Para dar un contraejemplo robusto en $p = \infty$, basta con elegir la matriz $J$ donde todos sus elementos son un $1$. Su norma vectorial infinita (el mayor elemento en valor absoluto) es $1$. Sin embargo, su norma de operador infinito (la máxima suma de los valores absolutos de una fila) es $N$. Como $N \neq 1$, la desigualdad también se sostiene para el caso infinito.
}


% Ejercicio 15
\ejercicio{resuelto}{profesor}{
    Suponga que $\mathbb{K} \in \{\mathbb{R}, \mathbb{C}\}$, $\alpha \in \mathbb{R} \cup \{-\infty\}$, $\beta \in \mathbb{R} \cup \{+\infty\}$, $\alpha < \beta$, y $a_{ij}, b_i \in C((\alpha, \beta); \mathbb{K}^N)$ para todo $i, j \in \{1, \dots, N\}$. Demuestre que, para cada $t_0 \in (\alpha, \beta)$ y $u_0 \in \mathbb{K}^N$, el problema
    $$\begin{cases} u' = A(t)u + B(t), \\ u(t_0) = u_0, \end{cases}$$
    donde $A(t) = (a_{ij}(t))_{1 \le i, j \le N}$ y $B(t) = (b_i(t))_{1 \le i \le N}$, posee una solución única en $(\alpha, \beta)$.
}{
    Dado que ya sabemos que el teorema de existencia y unicidad es válido para cualquier intervalo compacto (cerrado y acotado), vamos a construir el intervalo abierto $(\alpha, \beta)$ como una unión infinita de intervalos compactos anidados.

    Dado que $t_0 \in (\alpha, \beta)$, podemos construir una sucesión de intervalos compactos anidados $I_n = [\alpha_n, \beta_n]$ tales que:
    \begin{enumerate}
        \item $t_0 \in I_1$ (y por lo tanto, $t_0 \in I_n$ para todo $n$).
        \item $I_n \subset I_{n+1}$ para todo $n \in \mathbb{N}$.
        \item La unión de todos estos intervalos cubre todo el intervalo abierto: $\bigcup_{n=1}^\infty I_n = (\alpha, \beta)$.
    \end{enumerate}

    Para ser rigurosos, podemos definir $\alpha_n$ y $\beta_n$ de la siguiente manera:
    \begin{itemize}
        \item Si $\alpha > -\infty$, tomamos $\alpha_n = \alpha + \frac{t_0 - \alpha}{n+1}$. Si $\alpha = -\infty$, tomamos $\alpha_n = t_0 - n$.
        \item Si $\beta < +\infty$, tomamos $\beta_n = \beta - \frac{\beta - t_0}{n+1}$. Si $\beta = +\infty$, tomamos $\beta_n = t_0 + n$.
    \end{itemize}

    Por la hipótesis del problema, sabemos que el problema de valor inicial está resuelto para intervalos compactos con extremos reales. Como la matriz $A(t)$ y el vector $B(t)$ son continuos en $(\alpha, \beta)$, también lo son cuando los restringimos a cualquier subintervalo compacto $I_n \subset (\alpha, \beta)$.

    Por lo tanto, para cada $n \in \mathbb{N}$, existe una única función $u_n: I_n \to \mathbb{K}^N$ que es solución del problema:
    \[
    \begin{cases} u_n' = A(t)u_n + B(t), \quad t \in I_n \\ u_n(t_0) = u_0 \end{cases}
    \]

    Ahora debemos asegurarnos de que estas soluciones encajen entre sí. Supongamos que tomamos dos índices $m > n$. Esto significa que $I_n \subset I_m$. 

    La función $u_m$ es solución en $I_m$, por lo que su restricción al intervalo más pequeño $I_n$ (es decir, $u_m|_{I_n}$) sigue siendo una solución de la ecuación diferencial en $I_n$ que cumple la condición inicial $u_m(t_0) = u_0$. 

    Pero por la unicidad de soluciones en el intervalo compacto $I_n$ (nuestra hipótesis), cualquier solución en $I_n$ que cumpla la condición inicial debe ser exactamente $u_n$. Por lo tanto, concluimos que:
    \[
    u_m(t) = u_n(t) \quad \text{para todo } t \in I_n
    \]

    Gracias a la consistencia demostrada en el paso anterior, podemos definir una función global $u: (\alpha, \beta) \to \mathbb{K}^N$ bien definida. 

    Para cualquier $t \in (\alpha, \beta)$, existirá algún $N_0 \in \mathbb{N}$ tal que $t \in I_{N_0}$. Definimos:
    \[
    u(t) = u_{N_0}(t)
    \]

    Esta función está bien definida (no depende del intervalo elegido) porque si $t$ también pertenece a otro intervalo $I_m$ con $m > N_0$, ya demostramos que $u_m(t) = u_{N_0}(t)$. 

    Como para cualquier $t \in (\alpha, \beta)$, $u(t)$ coincide en un entorno de $t$ con una función derivable que satisface la ecuación diferencial, se deduce directamente que $u(t)$ es derivable en $(\alpha, \beta)$ y cumple:
    \[
    u'(t) = A(t)u(t) + B(t) \quad \text{para todo } t \in (\alpha, \beta)
    \]
    Además, como $t_0 \in I_1$, $u(t_0) = u_1(t_0) = u_0$. Por tanto,  existe una solución en todo $(\alpha, \beta)$.

    Supongamos que existe otra solución $v: (\alpha, \beta) \to \mathbb{K}^N$ para el mismo problema de valor inicial. 

    Para cualquier $t \in (\alpha, \beta)$, sabemos que existe un $n \in \mathbb{N}$ tal que $t \in I_n$. Si restringimos la función $v$ al intervalo compacto $I_n$, obtenemos una solución al problema en $I_n$. 
    Por la hipótesis de unicidad en intervalos compactos, la restricción de $v$ a $I_n$ debe ser exactamente $u_n$. Es decir:
    \[
    v(t) = u_n(t) = u(t)
    \]
    Como esto es válido para cualquier $t \in (\alpha, \beta)$, concluimos que $v(t) = u(t)$ en todo el dominio. Esto demuestra que la solución global es única.
}


% Ejercicio 16
\ejercicio{resuelto}{profesor}{
    Para cada uno de los siguientes problemas de valor inicial, calcule las iteraciones de Picard-Lindelöf asociadas y demuestre que convergen a la solución única correspondiente:
    \begin{multicols}{2}
        \begin{enumerate}[label=\alph*)]
            \item $\begin{cases} u' = u - 1, \\ u(0) = u_0 \in \mathbb{R}. \end{cases}$
            \item $\begin{cases} u'' = u, \\ u(0) = 1, \quad u'(0) = -1. \end{cases}$
            \item $\begin{cases} u'' = -u, \\ u(0) = 1, \quad u'(0) = 0. \end{cases}$
            \item $\begin{cases} u'' = -u, \\ u(0) = 0, \quad u'(0) = 1. \end{cases}$
        \end{enumerate}
    \end{multicols}
}{
    El método de iteraciones de Picard-Lindelöf nos permite construir una sucesión de funciones que converge a la solución exacta de un problema de valor inicial $u' = f(t, u)$ con $u(0) = u_0$. La fórmula recursiva general es:
    $$u_{n+1}(t) = u_0 + \int_0^t f(s, u_n(s)) \, ds$$

    Para las ecuaciones de segundo orden (incisos b, c y d), el procedimiento estándar es transformarlas en un sistema de primer orden $U' = AU$, donde $U(t) = \begin{pmatrix} u(t) \\ u'(t) \end{pmatrix}$. La iteración se aplica entonces de forma vectorial:
    $$U_{n+1}(t) = U_0 + \int_0^t A U_n(s) \, ds$$

    \begin{enumerate} [label=\alph*)]
        \item  $u' = u - 1, \quad u(0) = u_0 \in \mathbb{R}$\\
            En este caso, nuestra función es $f(t, u) = u - 1$.
            \begin{itemize}
                \item \textbf{Iteración 0:} Arrancamos con la condición inicial constante.
                $$u_0(t) = u_0$$
                \item \textbf{Iteración 1:}
                $$u_1(t) = u_0 + \int_0^t (u_0 - 1) \, ds = u_0 + (u_0 - 1)t$$
                \item \textbf{Iteración 2:}
                $$u_2(t) = u_0 + \int_0^t (u_1(s) - 1) \, ds = u_0 + \int_0^t \left(u_0 - 1 + (u_0 - 1)s\right) \, ds$$
                $$u_2(t) = u_0 + (u_0 - 1)t + (u_0 - 1)\frac{t^2}{2!}$$
                \item \textbf{Patrón general:} Si observamos atentamente y reescribimos el primer término como $u_0 = 1 + (u_0 - 1)$, el patrón para la $n$-ésima iteración es:
                $$u_n(t) = 1 + (u_0 - 1) \sum_{k=0}^n \frac{t^k}{k!}$$
                \item \textbf{Convergencia:} Tomando el límite cuando $n \to \infty$, reconocemos la serie de Taylor de la función exponencial $e^t$:
                $$u(t) = \lim_{n \to \infty} u_n(t) = 1 + (u_0 - 1)e^t$$
            \end{itemize}  
        \item $u'' = u, \quad u(0) = 1, \ u'(0) = -1$ \\
            Definimos el vector $U = \begin{pmatrix} u \\ v \end{pmatrix}$ donde $v = u'$. El sistema equivalente es $U' = \begin{pmatrix} 0 & 1 \\ 1 & 0 \end{pmatrix} U$, con condición inicial $U_0 = \begin{pmatrix} 1 \\ -1 \end{pmatrix}$.

            \begin{itemize}
                \item \textbf{Iteración 1:}
                $$U_1(t) = \begin{pmatrix} 1 \\ -1 \end{pmatrix} + \int_0^t \begin{pmatrix} 0 & 1 \\ 1 & 0 \end{pmatrix} \begin{pmatrix} 1 \\ -1 \end{pmatrix} ds = \begin{pmatrix} 1 \\ -1 \end{pmatrix} + \int_0^t \begin{pmatrix} -1 \\ 1 \end{pmatrix} ds = \begin{pmatrix} 1 - t \\ -1 + t \end{pmatrix}$$
                \item \textbf{Iteración 2:}
                $$U_2(t) = \begin{pmatrix} 1 \\ -1 \end{pmatrix} + \int_0^t \begin{pmatrix} 0 & 1 \\ 1 & 0 \end{pmatrix} \begin{pmatrix} 1 - s \\ -1 + s \end{pmatrix} ds = \begin{pmatrix} 1 \\ -1 \end{pmatrix} + \int_0^t \begin{pmatrix} -1 + s \\ 1 - s \end{pmatrix} ds = \begin{pmatrix} 1 - t + \frac{t^2}{2!} \\ -1 + t - \frac{t^2}{2!} \end{pmatrix}$$
                \item \textbf{Convergencia:} Nos interesa la primera componente, $u(t)$. El patrón es claramente alternante en signos pero manteniendo todas las potencias:
                $$u_n(t) = \sum_{k=0}^n \frac{(-1)^k t^k}{k!} = \sum_{k=0}^n \frac{(-t)^k}{k!}$$
                En el límite, esto converge a la serie exponencial:
                $$u(t) = e^{-t}$$
            \end{itemize}  
        \item $u'' = -u, \quad u(0) = 1, \ u'(0) = 0$ \\
            El sistema equivalente es $U' = \begin{pmatrix} 0 & 1 \\ -1 & 0 \end{pmatrix} U$, con $U_0 = \begin{pmatrix} 1 \\ 0 \end{pmatrix}$.

            \begin{itemize}
                \item \textbf{Iteración 1:}
                $$U_1(t) = U_0 + \int_0^t \begin{pmatrix} 0 & 1 \\ -1 & 0 \end{pmatrix} \begin{pmatrix} 1 \\ 0 \end{pmatrix} ds = \begin{pmatrix} 1 \\ 0 \end{pmatrix} + \int_0^t \begin{pmatrix} 0 \\ -1 \end{pmatrix} ds = \begin{pmatrix} 1 \\ -t \end{pmatrix}$$
                \item \textbf{Iteración 2:}
                $$U_2(t) = U_0 + \int_0^t \begin{pmatrix} 0 & 1 \\ -1 & 0 \end{pmatrix} \begin{pmatrix} 1 \\ -s \end{pmatrix} ds = \begin{pmatrix} 1 \\ 0 \end{pmatrix} + \int_0^t \begin{pmatrix} -s \\ -1 \end{pmatrix} ds = \begin{pmatrix} 1 - \frac{t^2}{2!} \\ -t \end{pmatrix}$$
                \item \textbf{Iteración 3:}
                $$U_3(t) = U_0 + \int_0^t \begin{pmatrix} 0 & 1 \\ -1 & 0 \end{pmatrix} \begin{pmatrix} 1 - \frac{s^2}{2!} \\ -s \end{pmatrix} ds = \begin{pmatrix} 1 \\ 0 \end{pmatrix} + \int_0^t \begin{pmatrix} -s \\ -1 + \frac{s^2}{2!} \end{pmatrix} ds = \begin{pmatrix} 1 - \frac{t^2}{2!} \\ -t + \frac{t^3}{3!} \end{pmatrix}$$
                \item \textbf{Convergencia:} La primera componente $u(t)$ acumula únicamente potencias pares con signos alternos:
                $$u_n(t) = \sum_{k=0}^{\lfloor n/2 \rfloor} \frac{(-1)^k t^{2k}}{(2k)!}$$
                Este límite corresponde exactamente a la serie de Taylor del coseno:
                $$u(t) = \cos(t)$$
            \end{itemize}
        \item $u'' = -u, \quad u(0) = 0, \ u'(0) = 1$ \\
            La matriz del sistema es la misma que en el inciso anterior, $A = \begin{pmatrix} 0 & 1 \\ -1 & 0 \end{pmatrix}$, pero ahora $U_0 = \begin{pmatrix} 0 \\ 1 \end{pmatrix}$.
            \begin{itemize}
                \item \textbf{Iteración 1:}
                $$U_1(t) = U_0 + \int_0^t \begin{pmatrix} 0 & 1 \\ -1 & 0 \end{pmatrix} \begin{pmatrix} 0 \\ 1 \end{pmatrix} ds = \begin{pmatrix} 0 \\ 1 \end{pmatrix} + \int_0^t \begin{pmatrix} 1 \\ 0 \end{pmatrix} ds = \begin{pmatrix} t \\ 1 \end{pmatrix}$$
                \item \textbf{Iteración 2:}
                $$U_2(t) = U_0 + \int_0^t \begin{pmatrix} 0 & 1 \\ -1 & 0 \end{pmatrix} \begin{pmatrix} s \\ 1 \end{pmatrix} ds = \begin{pmatrix} 0 \\ 1 \end{pmatrix} + \int_0^t \begin{pmatrix} 1 \\ -s \end{pmatrix} ds = \begin{pmatrix} t \\ 1 - \frac{t^2}{2!} \end{pmatrix}$$
                \item \textbf{Iteración 3:}
                $$U_3(t) = U_0 + \int_0^t \begin{pmatrix} 0 & 1 \\ -1 & 0 \end{pmatrix} \begin{pmatrix} s \\ 1 - \frac{s^2}{2!} \end{pmatrix} ds = \begin{pmatrix} 0 \\ 1 \end{pmatrix} + \int_0^t \begin{pmatrix} 1 - \frac{s^2}{2!} \\ -s \end{pmatrix} ds = \begin{pmatrix} t - \frac{t^3}{3!} \\ 1 - \frac{t^2}{2!} \end{pmatrix}$$
                \item \textbf{Convergencia:} Ahora, la primera componente $u(t)$ acumula potencias impares con signos alternos:
                $$u_n(t) = \sum_{k=0}^{\lfloor (n-1)/2 \rfloor} \frac{(-1)^k t^{2k+1}}{(2k+1)!}$$
                Este límite es la definición en serie de Taylor de la función seno:
                $$u(t) = \sin(t)$$
            \end{itemize}                           
    \end{enumerate}
}


% Ejercicio 17
\ejercicio{enunciado}{no}{
    Use inducción sobre $N$ para demostrar la identidad (1.23).
}{

}


% Ejercicio 18
\ejercicio{resuelto}{profesor}{
    Encuentre la solución general de la ecuación diferencial
    $$u'' - 3u' + 2u = \pi + e^{\sqrt{2}t}$$
    utilizando cada uno de los siguientes métodos:
    \begin{enumerate}[label=\alph*)]
        \item variación de parámetros (coeficientes), como se explica en la Sección 1.8;
        \item reducción del orden de la ecuación, como se explica en la Sección 1.9;
        \item mediante el método de coeficientes indeterminados, como se explica en la Sección 1.10.
    \end{enumerate}
}{
    La ecuación diferencial lineal de segundo orden con coeficientes constantes a resolver es:
    $$u'' - 3u' + 2u = \pi + e^{\sqrt{2}t}$$

    Para los tres métodos, necesitamos primero la solución de la ecuación homogénea asociada: $u'' - 3u' + 2u = 0$.
    Su ecuación característica es $r^2 - 3r + 2 = 0$, cuyas raíces son $r_1 = 1$ y $r_2 = 2$.
    Por lo tanto, la solución homogénea (complementaria) es:
    $$u_h(t) = C_1 e^t + C_2 e^{2t}$$
    Nuestro conjunto fundamental de soluciones es $u_1(t) = e^t$ y $u_2(t) = e^{2t}$.

    \begin{enumerate} [label=\alph*)]
        \item \textbf{Método de Variación de Parámetros} \\
            Buscamos una solución particular de la forma $u_p(t) = v_1(t)u_1(t) + v_2(t)u_2(t)$.
            Primero, calculamos el Wronskiano de $u_1$ y $u_2$:
            $$|W| = \begin{vmatrix} e^t & e^{2t} \\ e^t & 2e^{2t} \end{vmatrix} = 2e^{3t} - e^{3t} = e^{3t}$$

            Usamos las fórmulas para $v_1'$ y $v_2'$, considerando que $g(t) = \pi + e^{\sqrt{2}t}$:
            $$v_1' = \frac{-u_2(t) g(t)}{W} = \frac{-e^{2t}(\pi + e^{\sqrt{2}t})}{e^{3t}} = -\pi e^{-t} - e^{(\sqrt{2}-1)t}$$
            $$v_2' = \frac{u_1(t) g(t)}{W} = \frac{e^t(\pi + e^{\sqrt{2}t})}{e^{3t}} = \pi e^{-2t} + e^{(\sqrt{2}-2)t}$$

            Integramos para hallar $v_1$ y $v_2$:
            $$v_1(t) = \int \left( -\pi e^{-t} - e^{(\sqrt{2}-1)t} \right) dt = \pi e^{-t} - \frac{1}{\sqrt{2}-1} e^{(\sqrt{2}-1)t}$$
            $$v_2(t) = \int \left( \pi e^{-2t} + e^{(\sqrt{2}-2)t} \right) dt = -\frac{\pi}{2} e^{-2t} + \frac{1}{\sqrt{2}-2} e^{(\sqrt{2}-2)t}$$

            Construimos la solución particular $u_p(t) = v_1 u_1 + v_2 u_2$:
            $$u_p(t) = \left( \pi e^{-t} - \frac{1}{\sqrt{2}-1} e^{(\sqrt{2}-1)t} \right) e^t + \left( -\frac{\pi}{2} e^{-2t} + \frac{1}{\sqrt{2}-2} e^{(\sqrt{2}-2)t} \right) e^{2t}$$
            $$u_p(t) = \pi - \frac{1}{\sqrt{2}-1} e^{\sqrt{2}t} - \frac{\pi}{2} + \frac{1}{\sqrt{2}-2} e^{\sqrt{2}t}$$
            Agrupando términos, obtenemos:
            $$u_p(t) = \frac{\pi}{2} + \left( \frac{1}{\sqrt{2}-2} - \frac{1}{\sqrt{2}-1} \right) e^{\sqrt{2}t}$$

            Simplificamos el coeficiente de la exponencial buscando denominador común:
            $$\frac{1}{\sqrt{2}-2} - \frac{1}{\sqrt{2}-1} = \frac{(\sqrt{2}-1) - (\sqrt{2}-2)}{(\sqrt{2}-2)(\sqrt{2}-1)} = \frac{1}{4 - 3\sqrt{2}}$$
            Por lo tanto, la solución general $u(t) = u_h(t) + u_p(t)$ es:
            $$u(t) = C_1 e^t + C_2 e^{2t} + \frac{\pi}{2} + \frac{1}{4 - 3\sqrt{2}} e^{\sqrt{2}t}$$  
        \item \textbf{Reducción del Orden} \\
            Conocemos una solución homogénea $u_1(t) = e^t$. Proponemos que la solución tiene la forma $u(t) = v(t) e^t$.
            Calculamos sus derivadas:
            $$u' = v'e^t + ve^t$$
            $$u'' = v''e^t + 2v'e^t + ve^t$$

            Sustituimos en la EDO original:
            $$(v''e^t + 2v'e^t + ve^t) - 3(v'e^t + ve^t) + 2ve^t = \pi + e^{\sqrt{2}t}$$
            Simplificando los términos con $v$ y dividiendo por $e^t$, la ecuación se reduce a:
            $$v'' - v' = \pi e^{-t} + e^{(\sqrt{2}-1)t}$$

            Hacemos el cambio de variable $w = v'$, obteniendo una EDO de primer orden lineal para $w$:
            $$w' - w = \pi e^{-t} + e^{(\sqrt{2}-1)t}$$
            El factor integrante es $\mu(t) = e^{\int -1 dt} = e^{-t}$. Multiplicamos:
            $$(w e^{-t})' = \pi e^{-2t} + e^{(\sqrt{2}-2)t}$$
            Integramos respecto a $t$:
            $$w e^{-t} = -\frac{\pi}{2} e^{-2t} + \frac{1}{\sqrt{2}-2} e^{(\sqrt{2}-2)t} + C_1$$
            Despejamos $w$ (es decir, $v'$):
            $$v' = -\frac{\pi}{2} e^{-t} + \frac{1}{\sqrt{2}-2} e^{(\sqrt{2}-1)t} + C_1 e^t$$
            Integramos nuevamente para encontrar $v(t)$:
            $$v(t) = \frac{\pi}{2} e^{-t} + \frac{1}{(\sqrt{2}-2)(\sqrt{2}-1)} e^{(\sqrt{2}-1)t} + C_1 e^t + C_2$$
            Sabiendo que el denominador $(\sqrt{2}-2)(\sqrt{2}-1) = 4 - 3\sqrt{2}$, y multiplicando todo por $e^t$ para recuperar $u(t) = v(t)e^t$, obtenemos:
            $$u(t) = \frac{\pi}{2} + \frac{1}{4 - 3\sqrt{2}} e^{\sqrt{2}t} + C_1 e^{2t} + C_2 e^t$$
            Que es exactamente el mismo resultado.        
        \item \textbf{Método de Coeficientes Indeterminados} \\
            Dado que $g(t) = \pi + e^{\sqrt{2}t}$, proponemos una solución particular con la forma de la función de excitación:
            $$u_p(t) = A + B e^{\sqrt{2}t}$$
            \textit{Nota: Como ni $0$ ni $\sqrt{2}$ son raíces de nuestra ecuación característica, no es necesario multiplicar la propuesta por $t$.}

            Calculamos las derivadas:
            $$u_p'(t) = \sqrt{2}B e^{\sqrt{2}t}$$
            $$u_p''(t) = 2B e^{\sqrt{2}t}$$

            Sustituimos en la EDO original $u'' - 3u' + 2u = \pi + e^{\sqrt{2}t}$:
            $$(2B e^{\sqrt{2}t}) - 3(\sqrt{2}B e^{\sqrt{2}t}) + 2(A + B e^{\sqrt{2}t}) = \pi + e^{\sqrt{2}t}$$
            Agrupamos los términos constantes y los que acompañan a la exponencial:
            $$2A + (2 - 3\sqrt{2} + 2)B e^{\sqrt{2}t} = \pi + e^{\sqrt{2}t}$$
            $$2A + (4 - 3\sqrt{2})B e^{\sqrt{2}t} = \pi + e^{\sqrt{2}t}$$

            Igualamos los coeficientes a ambos lados de la ecuación:
            \begin{itemize}
                \item Para la constante: $2A = \pi \implies A = \frac{\pi}{2}$
                \item Para la exponencial: $(4 - 3\sqrt{2})B = 1 \implies B = \frac{1}{4 - 3\sqrt{2}}$
            \end{itemize}

            Por lo tanto, la solución general $u(t) = u_h(t) + u_p(t)$ es:
            $$u(t) = C_1 e^t + C_2 e^{2t} + \frac{\pi}{2} + \frac{1}{4 - 3\sqrt{2}} e^{\sqrt{2}t}$$            
    \end{enumerate}
}

% Ejercicio 19
\ejercicio{enunciado}{no}{
    Repita el ejercicio anterior con la siguiente ecuación diferencial:
    $$u'' - 2u' + u = 1 + e^{2t}.$$
}{

}

% Ejercicio 20
\ejercicio{resuelto}{profesor}{
    Encuentre la solución única del problema
    $$\begin{cases} u'' + u = \sec^3 t, \\ u(0) = 1, \ u'(0) = \frac{1}{2} \end{cases}$$
    en el intervalo $\left(-\frac{\pi}{2}, \frac{\pi}{2}\right)$.
}{
    En primer lugar se obtiene rapidamente que dos soluciones del homogéneo asociado $h'' + h = 0$ son $h_1(t) = \cos t$ y $h_2(t) = \sin t$. Ahora podemos construir el sistema de 1er orden equivalente a la EDO de 2º orden:
    \[
        \begin{pmatrix}
            u \\ v
        \end{pmatrix}'
        =
        \begin{pmatrix}
            0 & 1 \\
            -1 & 0
        \end{pmatrix}
        \begin{pmatrix}
            u \\ v
        \end{pmatrix}
        +
        \begin{pmatrix}
            0 \\ \sec^3 t
        \end{pmatrix}
    \]
    El Wronskiano de $h_1$ y $h_2$ es:
    \[
        W =
        \begin{pmatrix}
            \cos t & \sin t \\
            -\sin t & \cos t
        \end{pmatrix}
        \implies 
        |W| = \cos^2 t + \sin^2 t = 1
    \]
    Ahora, aplicamos el cambio de variable $U(t) = W(t) C(t)$, de lo que podemos pasar:
    \[
        U' = AU + B \implies
        \underbrace{W C'}_{AWC} + W' C = A W C + B \implies
        W C' = B \implies
        \begin{pmatrix}
            \cos t & \sin t \\
            -\sin t & \cos t
        \end{pmatrix}
        \begin{pmatrix}
            C_1' \\ C_2'
        \end{pmatrix}
        =
        \begin{pmatrix}
            0 \\ \sec^3 t
        \end{pmatrix}
    \]
    De lo que podemos despejar ahora que:
    \[
        C_1' =
        \begin{vmatrix}
            0 & \sin t \\
            \sec^3 t & \cos t
        \end{vmatrix}
        = -\sin t \sec^3 t =
        - \tan t \sec^2 t =
        - \tan t \cdot (\tan t)' =
        - \frac{1}{2} (\tan^2 t)'
    \]
    \[
        \implies \boxed{C_1 = -\frac{1}{2} \tan^2 t + x}
    \]
    \[
        C_2' =
        \begin{vmatrix}
            \cos t & 0 \\
            -\sin t & \sec^3 t
        \end{vmatrix}
        = \cos t \sec^3 t =
        \sec^2 t =
        (\tan t)'
    \]
    \[
        \implies \boxed{C_2 = \tan t + y}
    \]
    Por lo tanto, como $U = W C$, la solución general de la EDO es:
    \[
        u(t) = \cos t \cdot \left(-\frac{1}{2} \tan^2 t + x\right) + \sin t \cdot (\tan t + y) =
        x \cos t + y \sin t - \frac{1}{2} \cos t \tan^2 t + \sin t \tan t =
    \]
    \[
        = x \cos t + y \sin t + \frac{1}{2} \sin t \tan t
    \]
    Aplicando ahora las condiciones iniciales:
    \[
        u(0) = x \cos(0) + y \sin(0) + \frac{1}{2} \sin(0) \tan(0) = x = 1 \implies 
        x = 1
    \]
    \[
        u'(0) = -x \sin(0) + y \cos(0) + \frac{1}{2} (\cos(0) \tan(0) + \sin(0) \sec^2(0)) = y - \frac{1}{2} = \frac{1}{2} \implies
        y = 1
    \]
    Por lo tanto, $x = 1$ y $y = 1$. La solución particular es:
    \[
        u(t) = \cos t + \sin t + \frac{1}{2} \sin t \tan t
    \]
}


% Ejercicio 21
\ejercicio{resuelto}{profesor}{
    Suponga que $\gamma$ y $\omega$ son números reales positivos (distintos) y que $F_0 \in \mathbb{R}$.
    \begin{enumerate}
        \item[(a)] Encuentre la solución, $u_\gamma(t)$, del problema
        $$\begin{cases} u'' + \omega^2 u = F_0 \cos(\gamma t), \\ u(0) = u'(0) = 0. \end{cases}$$
        \item[(b)] Encuentre la solución, $u_\omega(t)$, del problema
        $$\begin{cases} u'' + \omega^2 u = F_0 \cos(\omega t), \\ u(0) = u'(0) = 0. \end{cases}$$
        \item[(c)] Demuestre que
        $$u_\omega(t) = \lim_{\gamma \to \omega} u_\gamma(t), \quad t \in \mathbb{R}.$$
        \item[(d)] Suponga que $\gamma \neq \omega$, $\gamma \sim \omega$, y demuestre que $u_\gamma(t)$, determinado en el inciso (a), tiene un comportamiento oscilatorio con una amplitud que a su vez oscila en el tiempo, teniendo una frecuencia pequeña. Determine la frecuencia y represente la gráfica de la solución. En Física, este fenómeno se conoce como \textbf{pulsación} (beat).
        \item[(e)] Represente la gráfica de $u_\omega(t)$. El fenómeno físico subyacente se conoce como \textbf{resonancia}. Compare la gráfica de $u_\omega(t)$ con la de $u_\gamma(t)$, para $\gamma \neq \omega, \gamma \sim \omega$.
    \end{enumerate}
}{
    Veamos los distintos apartados:
    \begin{enumerate}
        \item[(a)] Primero obtengamos la solución homogénea asociada a la ecuación diferencial:
        \[
            h'' + \omega^2 h = 0 \implies
            h'' = -\omega^2 h \implies
            \begin{cases}
                h_1(t) = \cos(\omega t) \\
                h_2(t) = \sin(\omega t)
            \end{cases}
        \]
        Ahora podemos:
        \[
            u_p(t) = m \cos (\gamma t) = \frac{F_0}{\omega^2 - \gamma^2} \cos(\gamma t) \implies
            m \gamma^2 \cos(\gamma t) + m \omega^2 \cos(\gamma t) = F_0 \cos(\gamma t) \implies 
        \]
        \[
            - m (\omega^2 - \gamma^2) \cos(\gamma t) = F_0 \cos(\gamma t) \implies
            m = \frac{F_0}{\omega^2 - \gamma^2}
        \]
        Y si ahora variamos los parámetros, proponiendo $u_p(t) = x_1 \cos t + x_2 \sin t$, obtenemos:
        \[
            u_p = x_1 \cos t + x_2 \sin t + \frac{F_0}{\omega^2 - \gamma^2} \cos(\gamma t) \implies
            u(0) = x_1 + \frac{F_0}{\omega^2 - \gamma^2} = 0 \implies
            x_1 = -\frac{F_0}{\omega^2 - \gamma^2}
        \]
        \[
            u'(t) = -x_1 \sin t + x_2 \cos t - \frac{F_0 \gamma}{\omega^2 - \gamma^2} \sin(\gamma t) \implies
            u'(0) = x_2 = 0
        \]
        Y por lo tanto, la solución particular es:
        \[
            u_\gamma(t) = \frac{F_0}{\omega^2 - \gamma^2} (\cos(\gamma t) - \cos(\omega t))
        \]
        \item[(b)] Para este caso, reducimos la ecuacion a un sistema de orden reducido, donde $u'=v$ y $v' = u'' = -\omega^2 u + F_0 \cos(\omega t)$. El sistema resultante es:
        \[
            \begin{pmatrix}
                u \\ v
            \end{pmatrix}'
            =
            \begin{pmatrix}
                0 & 1 \\
                -\omega^2 & 0
            \end{pmatrix}
            \begin{pmatrix}
                u \\ v
            \end{pmatrix}
            +
            \begin{pmatrix}
                0 \\ F_0 \cos(\omega t)
            \end{pmatrix}
        \]
        El Wronskiano de las soluciones del homogéneo asociado es:
        \[
            W =
            \begin{pmatrix}
                \cos(\omega t) & \sin(\omega t) \\
                -\omega \sin(\omega t) & \omega \cos(\omega t)
            \end{pmatrix}
            \implies
            |W| = \omega
        \]
        Ahora, aplicamos el cambio de variable $U(t) = W(t) C(t)$, de lo que podemos pasar:
        \[
            \underbrace{W C'}_{AWC} + W' C = A W C + B \implies
            W C' = B \implies
            \begin{pmatrix}
                \cos(\omega t) & \sin(\omega t) \\
                -\omega \sin(\omega t) & \omega \cos(\omega t)
            \end{pmatrix}
            \begin{pmatrix}
                C_1' \\ C_2'
            \end{pmatrix}
            =
            \begin{pmatrix}
                0 \\ F_0 \cos(\omega t)
            \end{pmatrix}
        \]
        De lo que despejamos ahora que:
        \[
            C_1' =
            \frac{1}{\omega}
            \begin{vmatrix}
                0 & \sin(\omega t) \\
                F_0 \cos(\omega t) & \omega \cos(\omega t)
            \end{vmatrix}
            = - \frac{F_0}{\omega} \sin(\omega t) \cos(\omega t) =
            - \frac{F_0}{2\omega} \sin(2\omega t)
        \]
        \[
            \implies \boxed{C_1 = \frac{F_0}{4\omega^2} \cos(2\omega t) + x_1}
        \]
        \[
            C_2' =
            \frac{1}{\omega}
            \begin{vmatrix}
                \cos(\omega t) & 0 \\
                -\omega \sin(\omega t) & F_0 \cos(\omega t)
            \end{vmatrix}
            = \frac{F_0}{\omega} \cos^2(\omega t) =
            \frac{F_0}{2\omega} (1 + \cos(2\omega t))
        \]
        \[
            \implies \boxed{C_2 = \frac{F_0}{2\omega} t + \frac{F_0}{4\omega^2} \sin(2\omega t) + x_2}
        \]
        Notese que para $C_2$ tambien se podria haber integrado el coseno cuadrado mediante integración por partes, obteniendo el mismo resultado.\\
        Ahora, como $U = W C$, la solución general de la EDO es:
        \[
            u(t) = \cos(\omega t) \cdot \left(\frac{F_0}{4\omega^2} \cos(2\omega t) + x_1\right) + \sin(\omega t) \cdot \left(\frac{F_0}{2\omega} t + \frac{F_0}{4\omega^2} \sin(2\omega t) + x_2\right) \implies
        \]
        Ahora si despejamos:
        \[
            u(t) = \frac{F_0}{4\omega^2} \cos(\omega t) \cos(2\omega t) + x_1 \cos(\omega t) + \frac{F_0}{2\omega} t \sin(\omega t) + \frac{F_0}{4\omega^2} \sin(\omega t) \sin(2\omega t) + x_2 \sin(\omega t)
        \]
        \[
            = x_1 \cos(\omega t) + x_2 \sin(\omega t) + \frac{F_0}{4\omega^2} \underbrace{ \left( \cos(\omega t) \cos(2\omega t) + \sin(\omega t) \sin(2\omega t) \right)}_{\cos(\omega t)} + \frac{F_0}{2\omega} t \sin(\omega t)
        \]
        Ahora imponemos las condiciones iniciales:
        \[
            u(t) = \left(x_1 + \frac{F_0}{4 \omega^2}\right) \cos(\omega t) + x_2 \sin(\omega t) + \frac{F_0}{2\omega} t \sin(\omega t)
        \]
        \[
            0 = u(0) = x_1 + \frac{F_0}{4 \omega^2} \implies x_1 = -\frac{F_0}{4 \omega^2}
        \]
        \[
            u'(t) = -\omega \left(x_1 + \frac{F_0}{4 \omega^2}\right) \sin(\omega t) + x_2 \omega \cos(\omega t) + \frac{F_0}{2\omega} \sin(\omega t) + \frac{F_0}{2} t \cos(\omega t)
        \]
        \[
            0 = u'(0) = x_2 \omega \implies x_2 = 0      
        \]

        Por lo tanto, la solución particular es:
        \[
            u_\omega(t) = \frac{F_0}{2\omega} t \sin(\omega t)
        \]
    
        \item [(c)] Para este apartado, simplemente calculamos el límite:
        \[
            \lim_{\gamma \to \omega} u_\gamma(t) = \lim_{\gamma \to \omega} \frac{F_0}{\omega^2 - \gamma^2} (\cos(\gamma t) - \cos(\omega t))
        \]
        Aplicando la regla de l'Hôpital, obtenemos:
        \[
            \lim_{\gamma \to \omega} u_\gamma(t) = \lim_{\gamma \to \omega} \frac{F_0 (-t \sin(\gamma t))}{-2\gamma} = \frac{F_0 t}{2\omega} \sin(\omega t) = u_\omega(t)
        \]
        
        \item [(d)]
        \[
            \cos(a+b) - \cos(a-b) = -2 \sin a \sin b, \quad \cos(a+b) + \cos(a-b) = 2 \cos a \cos b
        \]
        \[
            \cos(a-b) = \cos a \cos b + \sin a \sin b, \quad \cos(a+b) = \cos a \cos b - \sin a \sin b
        \]
        \[
            \cos(\alpha) - \cos(\beta) = 2 \sin\left(\frac{\alpha + \beta}{2}\right) \sin\left(\frac{\alpha - \beta}{2}\right)
        \]
        \[
            \frac{F_0}{(\omega-\gamma)(\omega+\gamma)} (\cos(\gamma t) - \cos(\omega t)) =
            \frac{F_0}{(\omega-\gamma)(\omega+\gamma)} \cdot 2 \sin\left(\frac{\gamma t + \omega t}{2}\right) \sin\left(\frac{\gamma t - \omega t}{2}\right) 
        \]
        \[
            \gamma = \omega + \epsilon, \quad \epsilon \to 0 \implies
        \]
        \[
            = \frac{F_0}{\omega+\gamma} \cdot \frac{\sin\left(\frac{\gamma t + \omega t}{2}\right)}{\frac{\gamma + \omega}{2}} \cdot \frac{\sin\left(\frac{\gamma t - \omega t}{2}\right)}{\frac{\gamma - \omega}{2}} = \frac{F_0}{2\omega+\varepsilon} \frac{2}{\varepsilon} \sin\left(\frac{(2\omega + \varepsilon) t}{2}\right) \sin\left(\frac{\varepsilon t}{2}\right) = A(t) \sin\left((\omega + \frac{\varepsilon}{2}) t\right)
        \]
    \end{enumerate}
}

% Ejercicio 22
\ejercicio{enunciado}{no}{
    Use las soluciones determinadas en el ejercicio anterior para obtener las soluciones de la correspondiente ecuación diferencial con condiciones iniciales arbitrarias: $u(0) = u_0, u'(0) = v_0$.
}{

}

% Ejercicio 23
\ejercicio{resuelto}{no}{
    Sean $\lambda_j \in \mathbb{R}, j \in \{1, 2, 3\}$, tres números reales diferentes, y considere el operador diferencial
    $$P(D) := (D - \lambda_1)(D - \lambda_2)(D - \lambda_3).$$
    Encuentre la solución general de
    $$P(D)u = b(t),$$
    donde $b \in C([\alpha, \beta]; \mathbb{R})$ para algunos $\alpha < \beta$.
}{
    Para resolver esta ecuación diferencial, primero encontramos la solución general de la ecuación homogénea asociada $P(D)u = 0$. Dado que $P(D)$ se factoriza como $(D - \lambda_1)(D - \lambda_2)(D - \lambda_3)$, las soluciones del homogéneo son:
    \[
        u_h(t) = C_1 e^{\lambda_1 t} + C_2 e^{\lambda_2 t} + C_3 e^{\lambda_3 t}
    \]
    donde $C_1, C_2, C_3$ son constantes arbitrarias. Estas funciones forman un conjunto fundamental de soluciones para el homogéneo.

    Calculamos el Wronskiano de estas soluciones para asegurarnos de que son linealmente independientes:
    \[
        W =
        \begin{vmatrix}
            e^{\lambda_1 t} & e^{\lambda_2 t} & e^{\lambda_3 t} \\
            \lambda_1 e^{\lambda_1 t} & \lambda_2 e^{\lambda_2 t} & \lambda_3 e^{\lambda_3 t} \\
            \lambda_1^2 e^{\lambda_1 t} & \lambda_2^2 e^{\lambda_2 t} & \lambda_3^2 e^{\lambda_3 t}
        \end{vmatrix}
    \]
    Factorizando $e^{(\lambda_1 + \lambda_2 + \lambda_3) t}$, obtenemos:
    \[
        W = e^{(\lambda_1 + \lambda_2 + \lambda_3) t}
        \begin{vmatrix}
            1 & 1 & 1 \\
            \lambda_1 & \lambda_2 & \lambda_3 \\
            \lambda_1^2 & \lambda_2^2 & \lambda_3^2
        \end{vmatrix}
    \]
    El determinante de la matriz es el producto de las diferencias de los $\lambda_j$:
    \[
        \begin{vmatrix}
            1 & 1 & 1 \\
            \lambda_1 & \lambda_2 & \lambda_3 \\
            \lambda_1^2 & \lambda_2^2 & \lambda_3^2
        \end{vmatrix} = (\lambda_2 - \lambda_1)(\lambda_3 - \lambda_1)(\lambda_3 - \lambda_2) \neq 0
    \]
    Por lo tanto, las soluciones $e^{\lambda_1 t}, e^{\lambda_2 t}, e^{\lambda_3 t}$ son linealmente independientes, y el Wronskiano es:
    \[
        W = e^{(\lambda_1 + \lambda_2 + \lambda_3) t} (\lambda_2 - \lambda_1)(\lambda_3 - \lambda_1)(\lambda_3 - \lambda_2) \neq 0
    \]
    Aplicamos ahora el método de variación de coeficientes: $U(t) = W(t) C(t)$, 
    \[
        U'(t) = W'(t) C(t) + W(t) C'(t) = A W(t) C(t) + B(t)
    \]
    \[
        \implies W(t) C'(t) = B(t) \implies
        \begin{pmatrix}
            e^{\lambda_1 t} & e^{\lambda_2 t} & e^{\lambda_3 t} \\
            \lambda_1 e^{\lambda_1 t} & \lambda_2 e^{\lambda_2 t} & \lambda_3 e^{\lambda_3 t} \\
            \lambda_1^2 e^{\lambda_1 t} & \lambda_2^2 e^{\lambda_2 t} & \lambda_3^2 e^{\lambda_3 t}
        \end{pmatrix}
        \begin{pmatrix}
            C_1' \\ C_2' \\ C_3'
        \end{pmatrix}
        =
        \begin{pmatrix}
            0 \\ 0 \\ b(t)
        \end{pmatrix}
    \]
    Ahora, aplicando la regla de Cramer, obtenemos:
    \[
        C_1' = \frac{1}{W} \begin{vmatrix}
            0 & e^{\lambda_2 t} & e^{\lambda_3 t} \\
            0 & \lambda_2 e^{\lambda_2 t} & \lambda_3 e^{\lambda_3 t} \\
            b(t) & \lambda_2^2 e^{\lambda_2 t} & \lambda_3^2 e^{\lambda_3 t}
        \end{vmatrix} = \frac{b(t) \begin{vmatrix}
            e^{\lambda_2 t} & e^{\lambda_3 t} \\
            \lambda_2 e^{\lambda_2 t} & \lambda_3 e^{\lambda_3 t}
        \end{vmatrix} }{W} = \frac{b(t) e^{(\lambda_2 + \lambda_3) t} (\lambda_3 - \lambda_2)}{W}
    \]
    \[
        C_2' = \frac{1}{W} \begin{vmatrix}
            e^{\lambda_1 t} & 0 & e^{\lambda_3 t} \\
            \lambda_1 e^{\lambda_1 t} & 0 & \lambda_3 e^{\lambda_3 t} \\
            \lambda_1^2 e^{\lambda_1 t} & b(t) & \lambda_3^2 e^{\lambda_3 t}
        \end{vmatrix} = -\frac{b(t) e^{(\lambda_1 + \lambda_3) t} (\lambda_3 - \lambda_1)}{W}
    \]
    \[
        C_3' = \frac{1}{W} \begin{vmatrix}
            e^{\lambda_1 t} & e^{\lambda_2 t} & 0 \\
            \lambda_1 e^{\lambda_1 t} & \lambda_2 e^{\lambda_2 t} & 0 \\
            \lambda_1^2 e^{\lambda_1 t} & \lambda_2^2 e^{\lambda_2 t} & b(t)
        \end{vmatrix} = \frac{b(t) e^{(\lambda_1 + \lambda_2) t} (\lambda_2 - \lambda_1)}{W}
    \]
    Integrando $C_1', C_2', C_3'$, obtenemos
    \[
        C_1(t) = \int \frac{b(s) e^{(\lambda_2 + \lambda_3) s} (\lambda_3 - \lambda_2)}{W(s)} ds
    \]
    \[
        C_2(t) = -\int \frac{b(s) e^{(\lambda_1 + \lambda_3) s} (\lambda_3 - \lambda_1)}{W(s)} ds
    \]
    \[
        C_3(t) = \int \frac{b(s) e^{(\lambda_1 + \lambda_2) s} (\lambda_2 - \lambda_1)}{W(s)} ds
    \]
    Por lo tanto, la solución particular de la ecuación diferencial es:
    \[
        u_p(t) = C_1(t) e^{\lambda_1 t} + C_2(t) e^{\lambda_2 t} + C_3(t) e^{\lambda_3 t}
    \]

}

% Ejercicio 24
\ejercicio{resuelto}{profesor}{
    Sean $r, \kappa \in \mathbb{R}$. Demuestre que el cambio de escala temporal $t = e^s, s \in \mathbb{R}$, transforma la \textbf{ecuación de Euler}
    $$t^2 u'' + r t u' + \kappa u = 0, \quad t > 0,$$
    en una ecuación diferencial con coeficientes constantes. Encuentre una base para el conjunto de soluciones de la ecuación de Euler de acuerdo con todos los valores posibles de los parámetros $r$ y $\kappa$.
}{
    Primero definamos:
    \[
        v(s) = u(t) = u(e^s)
    \]
    Por lo tanto, derivando:
    \[
        v'(s) = u'(e^s) e^s = u'(t) t
    \]
    \[
        v''(s) = u''(e^s) e^{2s} + u'(e^s) e^s = u''(t) t^2 + \underbrace{u'(t) t}_{v'(s)}
    \]
    Fijemonos en que:
    \[
        t^2 u'' = v''(s) - v'(s) \, \qquad
        tu' = v'(s)
    \]
    Por lo tanto, sustituyendo en la ecuación de Euler, obtenemos:
    \[
        v''(s) - v'(s) + r v'(s) + \kappa v(s) = 0 \implies
        v''(s) + (r-1) v'(s) + \kappa v(s) = 0
    \]
    Resolviendo la ecuación característica de esta ecuación diferencial con coeficientes constantes, obtenemos:
    \[
        \lambda^2 + (r-1) \lambda + \kappa = 0 \implies
        \lambda = \frac{-(r-1) \pm \sqrt{(r-1)^2 - 4\kappa}}{2}
    \]
    Ahora podemos distinguir los siguientes casos:
    \begin{enumerate}
        \item Si $(r-1)^2 - 4\kappa > 0$, entonces las raíces son reales y distintas, y una base para el conjunto de soluciones de la ecuación de Euler es:
        \[
            v(s) = e^{\lambda_1 s} x_1 + e^{\lambda_2 s} x_2 \implies
            u(t) = t^{\lambda_1} x_1 + t^{\lambda_2} x_2
        \]
        \item Si $(r-1)^2 - 4\kappa = 0$, entonces las raíces son reales e iguales, y una base para el conjunto de soluciones de la ecuación de Euler es (usando que $t = e^s \iff s = \ln t$):
        \[
            v(s) = e^{\lambda s} x_1 + s e^{\lambda s} x_2 =
            e^{\lambda s} (x_1 + x_2 s) \implies
            u(t) = t^{\lambda} (x_1 + x_2 \ln t)
        \]
        \item Si $(r-1)^2 - 4\kappa < 0$, entonces las raíces son complejas conjugadas, y una base para el conjunto de soluciones de la ecuación de Euler es, deontando a las soluciones de la ecuación característica como $\lambda = \alpha + i \beta$ y $\overline{\lambda} = \alpha - i \beta$:
        \[
            v(s) = e^{\alpha s} (x_1 \cos(\beta s) + x_2 \sin(\beta s)) \implies
            u(t) = t^{\alpha} (x_1 \cos(\beta \ln t) + x_2 \sin(\beta \ln t))
        \]
    \end{enumerate}
}

% Ejercicio 25
\ejercicio{resuelto}{profesor}{
    Sean $\lambda, \mu, \omega \in \mathbb{C}$ tres números complejos distintos.
    \begin{enumerate}
        \item[(a)] Encuentre la solución única, $u_\omega$, del problema de Cauchy
        $$\begin{cases} (D - \lambda)(D - \mu)u = e^{\omega t}, \quad t \in \mathbb{R}, \\ u(0) = u'(0) = 0. \end{cases}$$
        \item[(b)] Encuentre la solución única, $u_\lambda$, de
        $$\begin{cases} (D - \lambda)(D - \mu)u = e^{\lambda t}, \quad t \in \mathbb{R}, \\ u(0) = u'(0) = 0. \end{cases}$$
        \item[(c)] Compare $u_\lambda$ y $u_\omega$.
        \item[(d)] Encuentre la solución única de
        $$\begin{cases} (D - \lambda)^2 u = e^{\lambda t}, \quad t \in \mathbb{R}, \\ u(0) = u'(0) = 0. \end{cases}$$
    \end{enumerate}
}{
    \begin{enumerate}
        \item[(a)] El polinomio característico asociado a la parte homogénea es $P(w) = (w-\lambda)(w-\mu)$. Puesto que $\omega \neq \lambda$ y $\omega \neq \mu$, tenemos que $P(\omega) \neq 0$. 
        La solución general es la suma de la homogénea y una particular propuesta de la forma $\frac{1}{P(\omega)} e^{\omega t}$:
        \[
            u(t) = A e^{\lambda t} + B e^{\mu t} + \frac{1}{P(\omega)} e^{\omega t}
        \]
        Imponemos las condiciones iniciales $u(0) = u'(0) = 0$:
        \[
            \begin{cases}
                A + B + \frac{1}{P(\omega)} = 0 \\
                A\lambda + B\mu + \frac{\omega}{P(\omega)} = 0
            \end{cases}
            \implies
            \begin{pmatrix} 1 & 1 \\ \lambda & \mu \end{pmatrix} \begin{pmatrix} A \\ B \end{pmatrix} = \frac{-1}{P(\omega)} \begin{pmatrix} 1 \\ \omega \end{pmatrix}
        \]
        Aplicando la regla de Cramer para despejar los coeficientes:
        \begin{align*}
            A &= \frac{\begin{vmatrix} -1/P(\omega) & 1 \\ -\omega/P(\omega) & \mu \end{vmatrix}}{\begin{vmatrix} 1 & 1 \\ \lambda & \mu \end{vmatrix}} = \frac{\frac{-\mu + \omega}{P(\omega)}}{\mu - \lambda} = \frac{\omega - \mu}{(\omega-\lambda)(\omega-\mu)(\mu-\lambda)} = \frac{1}{(\omega-\lambda)(\mu-\lambda)} \\
            B &= \frac{\begin{vmatrix} 1 & -1/P(\omega) \\ \lambda & -\omega/P(\omega) \end{vmatrix}}{\mu - \lambda} = \frac{\frac{-\omega + \lambda}{P(\omega)}}{\mu - \lambda} = \frac{-(\omega-\lambda)}{(\omega-\lambda)(\omega-\mu)(\mu-\lambda)} = \frac{1}{(\omega-\mu)(\lambda-\mu)}
        \end{align*}
        Sustituyendo $A$ y $B$, la solución única es:
        \[
            u_\omega(t) = \frac{e^{\lambda t}}{(\omega-\lambda)(\mu-\lambda)} + \frac{e^{\mu t}}{(\omega-\mu)(\lambda-\mu)} + \frac{e^{\omega t}}{(\omega-\lambda)(\omega-\mu)}
        \]

        \item[(b)] Para $(D-\lambda)(D-\mu)u = e^{\lambda t}$, como $\lambda$ es raíz simple del polinomio característico, proponemos una solución particular de la forma $u_p(t) = C t e^{\lambda t}$. 
        Evaluamos el operador diferencial:
        \begin{align*}
            (D-\mu)(D-\lambda)(C t e^{\lambda t}) &= (D-\mu) [C e^{\lambda t} + C \lambda t e^{\lambda t} - \lambda C t e^{\lambda t}] \\
            &= (D-\mu)(C e^{\lambda t}) = C \lambda e^{\lambda t} - C \mu e^{\lambda t} = C(\lambda-\mu)e^{\lambda t}
        \end{align*}
        Igualando a $e^{\lambda t}$, obtenemos $C = \frac{1}{\lambda-\mu}$.
        La solución general es: $u(t) = A_1 e^{\lambda t} + B_1 e^{\mu t} + \frac{t}{\lambda-\mu} e^{\lambda t}$.
        Imponiendo $u(0)=u'(0)=0$, el sistema resultante nos da $A_1 = -\frac{1}{(\lambda-\mu)^2}$ y $B_1 = \frac{1}{(\lambda-\mu)^2}$.
        Por tanto, la solución es:
        \[
            u_\lambda(t) = -\frac{e^{\lambda t}}{(\lambda-\mu)^2} + \frac{e^{\mu t}}{(\lambda-\mu)^2} + \frac{t e^{\lambda t}}{\lambda-\mu}
        \]

        \item[(c)] Comparamos tomando el límite cuando $\omega \to \lambda$ de la solución $u_\omega(t)$. Agrupamos los términos con denominadores que se anulan:
        \[
            \lim_{\omega \to \lambda} u_\omega(t) = \lim_{\omega \to \lambda} \left[ \frac{-e^{\lambda t}(\omega-\mu) + e^{\omega t}(\lambda-\mu)}{(\omega-\lambda)(\omega-\mu)(\lambda-\mu)} + \frac{e^{\mu t}}{(\omega-\mu)(\lambda-\mu)} \right]
        \]
        Aplicamos la regla de L'Hôpital derivando respecto a $\omega$ en el primer bloque:
        \[
            \lim_{\omega \to \lambda} \frac{-e^{\lambda t} + t e^{\omega t}(\lambda-\mu)}{(\lambda-\mu)(2\omega - \mu - \lambda)} = \frac{-e^{\lambda t} + t e^{\lambda t}(\lambda-\mu)}{(\lambda-\mu)(\lambda-\mu)} = -\frac{e^{\lambda t}}{(\lambda-\mu)^2} + \frac{t e^{\lambda t}}{\lambda-\mu}
        \]
        Sumando el término en $e^{\mu t}$ (cuyo límite se evalúa directamente sustituyendo $\omega$ por $\lambda$), recuperamos exactamente $u_\lambda(t)$. Luego queda demostrado que $u_\lambda(t) = \lim_{\omega \to \lambda} u_\omega(t)$.

        \item[(d)] Para $(D-\lambda)^2 u = e^{\lambda t}$, siendo $\lambda$ raíz doble del polinomio característico, proponemos la solución particular $u_p(t) = m t^2 e^{\lambda t}$.
        \begin{align*}
            (D-\lambda)^2 (m t^2 e^{\lambda t}) &= (D-\lambda) [2mt e^{\lambda t} + \lambda m t^2 e^{\lambda t} - \lambda m t^2 e^{\lambda t}] \\
            &= (D-\lambda)(2mt e^{\lambda t}) = 2m e^{\lambda t} + 2m\lambda t e^{\lambda t} - 2m\lambda t e^{\lambda t} = 2m e^{\lambda t}
        \end{align*}
        Igualando a $e^{\lambda t}$, resulta $2m = 1 \implies m = 1/2$.
        La solución general es $u(t) = A e^{\lambda t} + B t e^{\lambda t} + \frac{1}{2}t^2 e^{\lambda t}$.
        Con las condiciones $u(0)=u'(0)=0$, se deduce rápidamente que $A=0$ y $B=0$, por lo que la solución única es:
        \[
            u(t) = \frac{1}{2}t^2 e^{\lambda t}
        \]
    \end{enumerate}
}

% Ejercicio 26
\ejercicio{resuelto}{profesor}{
    26. Dado $L > 0$, sea $U_d$ el espacio vectorial de las funciones $u \in C^2([0, L]; \mathbb{R})$ tales que $u(0) = u(L) = 0$, y considere el operador diferencial
    $$-D^2 : U_d \to V := C([0, L]; \mathbb{R}),$$
    definido como
    $$-D^2 u = -u'' \text{ para todo } u \in U_d.$$
    Se dice que un número $\lambda \in \mathbb{R}$ es un \textbf{valor propio} (eigenvalue) del operador diferencial $-D^2$ en $U_d$ si existe $u \in U_d, u \neq 0$, tal que $-D^2 u = \lambda u$. En tal caso, se dice que $u$ es una \textbf{función propia} (eigenfunction) asociada al valor propio $\lambda$. Encuentre los valores propios y las funciones propias de $-D^2$ en $U_d$, y represente cuidadosamente todas las funciones propias. (El subíndice $d$ en $U_d$ se refiere a las condiciones $u(0) = u(L) = 0$, conocidas como condiciones de contorno de \textbf{Dirichlet}).
}{
    Buscamos soluciones no nulas para $-u'' = \lambda u \implies u'' + \lambda u = 0$, sometidas a las condiciones de frontera de Dirichlet $u(0) = 0$ y $u(L) = 0$. Analizamos los posibles valores de $\lambda \in \mathbb{R}$:

    \textbf{Caso 1: $\lambda < 0$.} Sea $\lambda = -\mu^2$ con $\mu > 0$.
    La ecuación es $u'' - \mu^2 u = 0$. La ecuación característica es $z^2 - \mu^2 = 0 \implies z = \pm \mu$.
    La solución general es $u(t) = x e^{\mu t} + y e^{-\mu t}$.
    Aplicamos las condiciones de frontera:
    \begin{align*}
        u(0) &= x + y = 0 \implies y = -x \\
        u(L) &= x e^{\mu L} + y e^{-\mu L} = x(e^{\mu L} - e^{-\mu L}) = 0
    \end{align*}
    Como $\mu > 0$ y $L > 0$, el término $(e^{\mu L} - e^{-\mu L})$ es estrictamente mayor que cero. Por lo tanto, necesariamente $x = 0$, lo que implica $y = 0$.
    La única solución es $u(t) \equiv 0$. Luego, el problema carece de autovalores negativos.

    \textbf{Caso 2: $\lambda = 0$.}
    La ecuación es $u'' = 0$. La ecuación característica tiene la raíz doble $z = 0$.
    La solución general es $u(t) = x + yt$.
    Aplicamos las condiciones de frontera:
    \begin{align*}
        u(0) &= x = 0 \\
        u(L) &= x + yL = yL = 0 \implies y = 0 \quad (\text{ya que } L > 0)
    \end{align*}
    Nuevamente obtenemos la solución trivial $u(t) \equiv 0$. Por tanto, $\lambda = 0$ no es autovalor.

    \textbf{Caso 3: $\lambda > 0$.}
    La ecuación es $u'' + \lambda u = 0$. La ecuación característica es $z^2 + \lambda = 0 \implies z = \pm i\sqrt{\lambda}$.
    La solución general es $u(t) = x \cos(\sqrt{\lambda}t) + y \sin(\sqrt{\lambda}t)$.
    Aplicamos las condiciones de frontera:
    \begin{align*}
        u(0) &= x \cos(0) + y \sin(0) = x = 0 \\
        u(L) &= y \sin(\sqrt{\lambda}L) = 0
    \end{align*}
    Para tener funciones propias (soluciones no nulas, es decir, $u \neq 0$), necesitamos que $y \neq 0$. Esto obliga a que el seno se anule:
    \[
        \sin(\sqrt{\lambda}L) = 0 \implies \sqrt{\lambda}L = k\pi, \quad \text{con } k = 1, 2, 3, \dots
    \]
    Despejando $\lambda$, encontramos la sucesión de valores propios:
    \[
        \lambda_k = \left( \frac{k\pi}{L} \right)^2, \quad k \ge 1
    \]
    Y las funciones propias asociadas (tomando la constante libre $y=1$) son:
    \[
        u_k(t) = \sin\left(\frac{k\pi}{L}t\right)
    \]
}

% Ejercicio 27
\ejercicio{enunciado}{no}{
    Dado $L > 0$, sea $U_n$ el espacio vectorial de las funciones $u \in C^2([0, L]; \mathbb{R})$ tales que $u'(0) = u'(L) = 0$, y considere el operador diferencial
    $$-D^2 : U_n \to V := C([0, L]; \mathbb{R}),$$
    definido como
    $$-D^2 u = -u'' \text{ para todo } u \in U_n.$$
    Se dice que un número $\lambda \in \mathbb{R}$ es un \textbf{valor propio} del operador diferencial $-D^2$ en $U_n$ si existe $u \in U_n, u \neq 0$, tal que $-D^2 u = \lambda u$. En tal caso, se dice que $u$ es una función propia asociada al valor propio $\lambda$. Encuentre los valores propios y las funciones propias de $-D^2$ en $U_n$, y represente cuidadosamente todas las funciones propias. (El subíndice $n$ en $U_n$ se refiere a las condiciones $u'(0) = u'(L) = 0$, conocidas como condiciones de contorno de \textbf{Neumann}).
}{

}

% Ejercicio 28
\ejercicio{enunciado}{no}{
    Demuestre que la función $u(t) = e^t$ resuelve la ecuación diferencial
    $$u'' - tu' + (t - 1)u = 0, \quad t \in \mathbb{R}.$$
    Encuentre la solución general de
    $$u'' - tu' + (t - 1)u = e^{\frac{t^2}{2}}.$$
}{

}

% Ejercicio 29
\ejercicio{resuelto}{profesor}{
    Dada la ecuación diferencial
    $$u'' + \frac{1}{t}u' + \left(1 - \frac{1}{4t^2}\right)u = 0, \quad t > 0, \quad (1.96)$$
    utilice el cambio de variable $u(t) = t^\alpha v(t)$, para una elección adecuada de $\alpha$, con el fin de encontrar la solución general de (1.96). Luego, encuentre la solución general de la ecuación no homogénea
    $$u'' + \frac{1}{t}u' + \left(1 - \frac{1}{4t^2}\right)u = 3\sqrt{t} \sin t, \quad t > 0.$$
}{
    Realizamos el cambio de variable propuesto $u(t) = t^\alpha v(t)$. Calculamos sus primera y segunda derivadas respecto a $t$:
    \begin{align*}
        u'(t) &= \alpha t^{\alpha-1}v(t) + t^\alpha v'(t) \\
        u''(t) &= \alpha(\alpha-1)t^{\alpha-2}v(t) + 2\alpha t^{\alpha-1}v'(t) + t^\alpha v''(t)
    \end{align*}
    Sustituimos estas expresiones en la ecuación diferencial homogénea (1.96):
    \begin{align*}
        & \left[ t^\alpha v''(t) + 2\alpha t^{\alpha-1}v'(t) + \alpha(\alpha-1)t^{\alpha-2}v(t) \right] \\
        + & \frac{1}{t} \left[ \alpha t^{\alpha-1}v(t) + t^\alpha v'(t) \right] + \left(1 - \frac{1}{4t^2}\right) t^\alpha v(t) = 0
    \end{align*}
    Agrupamos los términos según las derivadas de la nueva función incógnita $v$:
    \[
        t^\alpha v''(t) + (2\alpha t^{\alpha-1} + t^{\alpha-1}) v'(t) + \left( \alpha(\alpha-1)t^{\alpha-2} + \alpha t^{\alpha-2} + t^\alpha - \frac{1}{4}t^{\alpha-2} \right) v(t) = 0
    \]
    Para simplificar la ecuación resultante, buscamos un valor de $\alpha$ que anule por completo el coeficiente del término de primera derivada $v'(t)$:
    \[
        (2\alpha + 1)t^{\alpha-1} = 0 \implies 2\alpha + 1 = 0 \implies \alpha = -\frac{1}{2}
    \]
    Sustituyendo $\alpha = -1/2$ en la ecuación agrupada:
    \begin{align*}
        t^{-1/2} v''(t) + \left( \left(-\frac{1}{2}\right)\left(-\frac{3}{2}\right)t^{-5/2} - \frac{1}{2}t^{-5/2} + t^{-1/2} - \frac{1}{4}t^{-5/2} \right) v(t) &= 0 \\
        t^{-1/2} v''(t) + \left( \frac{3}{4}t^{-5/2} - \frac{3}{4}t^{-5/2} + t^{-1/2} \right) v(t) &= 0 \\
        t^{-1/2} v''(t) + t^{-1/2} v(t) &= 0
    \end{align*}
    Multiplicando toda la ecuación por $t^{1/2}$ (que es válido ya que $t > 0$), obtenemos la ecuación diferencial de un oscilador armónico simple:
    \[
        v''(t) + v(t) = 0 \implies v_h(t) = x_1 \sin t + x_2 \cos t
    \]
    Deshaciendo el cambio de variable original ($u = t^{-1/2}v$), obtenemos la solución general de la ecuación homogénea de Bessel:
    \[
        u_h(t) = \frac{x_1 \sin t + x_2 \cos t}{\sqrt{t}}
    \]

    A continuación, resolvemos la ecuación no homogénea. Al aplicar el mismo cambio de variable $u(t) = t^{-1/2} v(t)$, sabemos que el operador lineal del lado izquierdo se reduce exactamente a $t^{-1/2}(v''(t) + v(t))$. Lo igualamos al término forzante:
    \[
        \frac{1}{\sqrt{t}} \big( v''(t) + v(t) \big) = 3\sqrt{t} \sin t \implies v''(t) + v(t) = 3t \sin t
    \]
    Para resolver esta nueva ecuación con coeficientes constantes para $v$, proponemos una solución particular por coeficientes indeterminados debido a la resonancia con el término homogéneo:
    \[
        v_p(t) = (A t^2 + B t) \cos t + (C t^2 + D t) \sin t
    \]
    Derivando y ajustando los coeficientes (como se mostró en clase), se obtiene la solución para $v(t)$. Finalmente, la solución general buscada será:
    \[
        u(t) = \frac{v_h(t) + v_p(t)}{\sqrt{t}}
    \]
}

% Ejercicio 30
\ejercicio{enunciado}{no}{
    Encuentre la solución general de la ecuación diferencial no homogénea
    $$u''' + 3u'' + 3u' + u = \pi + e^t.$$
}{
    Resolvamos la ecuación diferencial dada:
    \[
        u''' + 3u'' + 3u' + u = \pi + e^t
    \]
    \textbf{1. Solución de la ecuación homogénea ($u_h$)}

    El polinomio característico asociado a la ecuación diferencial es:
    \[
        P(z) = z^3 + 3z^2 + 3z + 1
    \]
    Reconocemos de inmediato que esta expresión es el desarrollo de un binomio al cubo:
    \[
        P(z) = (z+1)^3
    \]
    Igualando a cero, vemos que la única raíz es $z = -1$, la cual tiene una multiplicidad de 3 (raíz triple). Por lo tanto, una base de soluciones para el espacio homogéneo $\Sigma_0$ es:
    \[
        \mathcal{B} = \{ e^{-t}, t e^{-t}, t^2 e^{-t} \}
    \]
    La solución general de la ecuación homogénea es:
    \[
        u_h(t) = c_1 e^{-t} + c_2 t e^{-t} + c_3 t^2 e^{-t}
    \]

    \textbf{2. Solución particular ($u_p$)}

    Por el principio de superposición, calculamos una solución particular para cada término del lado derecho $f(t) = \pi + e^t$.

    \textit{Para el término constante $\pi$:}
    Como $0$ no es raíz del polinomio característico ($P(0) = 1 \neq 0$), no hay resonancia. Proponemos una solución constante $u_{p1} = A$:
    \[
        P(D) u_{p1} = \pi \implies P(D) A = \pi \implies A = \pi
    \]
    (Al evaluar el operador diferencial sobre una constante, todas las derivadas se anulan y solo sobrevive el término $1 \cdot A$). Así, $u_{p1} = \pi$.

    \textit{Para el término exponencial $e^t$:}
    Como $1$ no es raíz del polinomio característico ($P(1) = (1+1)^3 = 8 \neq 0$), proponemos la forma estándar $u_{p2} = B e^t$.
    Utilizando la propiedad de las exponenciales $P(D)e^{\alpha t} = P(\alpha)e^{\alpha t}$, tenemos:
    \[
        P(D)(B e^t) = e^t \implies B \cdot P(1) e^t = e^t \implies 8B e^t = e^t \implies 8B = 1 \implies B = \frac{1}{8}
    \]
    Por lo tanto, $u_{p2} = \frac{1}{8}e^t$.

    \textbf{3. Solución general}

    La solución general de la ecuación diferencial es la suma de la solución homogénea y las soluciones particulares calculadas ($u = u_h + u_{p1} + u_{p2}$):
    \[
        u(t) = c_1 e^{-t} + c_2 t e^{-t} + c_3 t^2 e^{-t} + \pi + \frac{1}{8} e^t
    \]
    donde $c_1, c_2, c_3 \in \mathbb{R}$ son constantes arbitrarias.
}

% Ejercicio 31
\ejercicio{enunciado}{no}{
    Para un $\omega > 0$ dado, encuentre la solución general (real y compleja) de cada una de las siguientes ecuaciones diferenciales:
    \begin{enumerate}
        \item[(a)] $u^{(6)} - 8u^{(5)} + 25u^{(4)} - 32u''' - u'' + 40u' - 25u = 1 + e^{2t}$.
        \item[(b)] $u^{(4)} - 4u''' + 12u'' + 4u' - 13u = 1 + e^{2t} + \cos(\omega t)$.
        \item[(c)] $u^{(5)} + 3u^{(4)} + 3u''' + u'' = 1 + e^{2t}$.
    \end{enumerate}
}{
    \begin{enumerate}
        \item[(a)] 
        Resolvamos la ecuación diferencial:
        \[
            u^{(6)} - 8u^{(5)} + 25u^{(4)} - 32u''' - u'' + 40u' - 25u = 1 + e^{2t}
        \]

        \textbf{1. Solución de la ecuación homogénea ($u_h$)}

        El polinomio característico asociado a la ecuación homogénea es:
        \[
            P(z) = z^6 - 8z^5 + 25z^4 - 32z^3 - z^2 + 40z - 25
        \]
        Aplicando la regla de Ruffini, podemos comprobar que $1$ y $-1$ son raíces evidentes del polinomio. Al factorizar, obtenemos:
        \[
            P(z) = (z-1)(z+1)(z^4 - 8z^3 + 26z^2 - 40z + 25) 
        \]
        Para factorizar el polinomio de cuarto grado restante, nos apoyamos en el desarrollo de $(z-2)^4$:
        \[
            (z-2)^4 = z^4 - 8z^3 + 24z^2 - 32z + 16
        \]
        Reescribimos el factor de grado 4 buscando esta estructura:
        \begin{align*}
            z^4 - 8z^3 + 26z^2 - 40z + 25 &= (z-2)^4 + 2z^2 - 8z + 9 \\
            &= (z-2)^4 + 2(z^2 - 4z + 4) + 1 \\
            &= (z-2)^4 + 2(z-2)^2 + 1 \\
            &= \left( (z-2)^2 + 1 \right)^2
        \end{align*}
        Por lo tanto, la factorización completa del polinomio es:
        \[
            P(z) = (z-1)(z+1) \left( (z-2)^2 + 1 \right)^2
        \]
        De aquí deducimos que las raíces del polinomio característico son $1, -1$ (raíces reales simples) y $2 \pm i$ (raíces complejas conjugadas dobles). 

        En consecuencia, una base de soluciones reales para el espacio nulo $\Sigma_0$ es:
        \[
            \mathcal{B}_{\mathbb{R}} = \{ e^t, e^{-t}, e^{2t} \cos t, e^{2t} \sin t, t e^{2t} \cos t, t e^{2t} \sin t \}
        \]
        Y la base de soluciones complejas es:
        \[
            \mathcal{B}_{\mathbb{C}} = \{ e^t, e^{-t}, e^{(2+i)t}, e^{(2-i)t}, t e^{(2+i)t}, t e^{(2-i)t} \}
        \]

        \textbf{2. Solución particular ($u_p$)}

        Aplicamos el principio de superposición para el término no homogéneo $f(t) = 1 + e^{2t}$.

        Para el término constante $1$, como $0$ no es raíz del polinomio característico, proponemos una constante $u_{p1} = K$:
        \[
            P(D) u_{p1} = 1 \implies P(D) K = 1 \iff -25K = 1 \implies K = -\frac{1}{25}
        \]

        Para el término exponencial $e^{2t}$, como $2$ no es raíz de $P(z)$, proponemos $u_{p2} = m e^{2t}$. Utilizamos la propiedad $P(D)e^{\alpha t} = P(\alpha)e^{\alpha t}$:
        \[
            P(D) (m e^{2t}) = e^{2t} \implies m P(2) e^{2t} = e^{2t}
        \]
        Evaluamos $P(z)$ en $z=2$:
        \[
            P(2) = (2-1)(2+1)\left( (2-2)^2 + 1 \right)^2 = (1)(3)(1)^2 = 3
        \]
        Sustituyendo en la ecuación:
        \[
            m \cdot 3 \cdot e^{2t} = e^{2t} \implies 3m = 1 \implies m = \frac{1}{3}
        \]
        La solución particular total es:
        \[
            u_p(t) = -\frac{1}{25} + \frac{1}{3}e^{2t}
        \]

        \textbf{3. Solución general}

        La solución general de la ecuación diferencial es la suma de la solución homogénea y la particular ($u = u_h + u_p$).
        \[
            u(t) = c_1 e^t + c_2 e^{-t} + e^{2t}(c_3 \cos t + c_4 \sin t + c_5 t \cos t + c_6 t \sin t) - \frac{1}{25} + \frac{1}{3}e^{2t}
        \]
        donde $c_i \in \mathbb{R}$ son constantes arbitrarias.
        \item[(b)]
        Resolvamos la ecuación diferencial:
        \[
            u^{(4)} - 4u''' + 12u'' + 4u' - 13u = 1 + e^{2t} + \cos(\omega t) \quad (\text{con } \omega > 0)
        \]

        \textbf{1. Solución de la ecuación homogénea ($u_h$)}

        El polinomio característico asociado es:
        \[
            P(z) = z^4 - 4z^3 + 12z^2 + 4z - 13
        \]
        Buscamos raíces racionales probando los divisores del término independiente ($-13$). Evaluando, encontramos que $P(1) = 0$ y $P(-1) = 0$. Por tanto, el polinomio es divisible por $(z-1)(z+1) = z^2 - 1$.
        Dividiendo $P(z)$ entre $z^2 - 1$ obtenemos la factorización:
        \[
            P(z) = (z^2 - 1)(z^2 - 4z + 13)
        \]
        Las raíces del factor cuadrático se obtienen mediante la fórmula general:
        \[
            z = \frac{4 \pm \sqrt{16 - 52}}{2} = \frac{4 \pm \sqrt{-36}}{2} = 2 \pm 3i
        \]
        Las raíces son $1, -1$ (reales simples) y $2 \pm 3i$ (complejas conjugadas simples). La solución general de la ecuación homogénea es:
        \[
            u_h(t) = c_1 e^t + c_2 e^{-t} + e^{2t}(c_3 \cos(3t) + c_4 \sin(3t))
        \]

        \textbf{2. Solución particular ($u_p$)}

        Por el principio de superposición, buscamos $u_p = u_{p1} + u_{p2} + u_{p3}$ para cada término del lado derecho $f(t) = 1 + e^{2t} + \cos(\omega t)$.

        \textit{Para el término constante $1$:}
        Como $0$ no es raíz de $P(z)$ ($P(0) = -13$), proponemos $u_{p1} = K$:
        \[
            P(D)K = 1 \implies -13K = 1 \implies u_{p1} = -\frac{1}{13}
        \]

        \textit{Para el término exponencial $e^{2t}$:}
        Como $2$ no es raíz de $P(z)$, calculamos $P(2) = 16 - 32 + 48 + 8 - 13 = 27$. Proponemos $u_{p2} = m e^{2t}$:
        \[
            P(D)(m e^{2t}) = e^{2t} \implies m P(2) e^{2t} = e^{2t} \implies 27m = 1 \implies u_{p2} = \frac{1}{27}e^{2t}
        \]

        \textit{Para el término trigonométrico $\cos(\omega t)$:}
        Reescribimos el coseno usando la fórmula de Euler: $\cos(\omega t) = \text{Re}(e^{i\omega t})$. 
        Resolvemos primero la ecuación compleja $P(D)U = e^{i\omega t}$. Puesto que las raíces complejas de $P(z)$ tienen parte real $2$, el valor puramente imaginario $i\omega$ nunca será raíz del polinomio para ningún $\omega > 0$. Por tanto, no hay resonancia.
        Proponemos la solución compleja $U_p = M e^{i\omega t}$:
        \[
            P(D)(M e^{i\omega t}) = e^{i\omega t} \implies M \cdot P(i\omega) = 1 \implies U_p = \frac{e^{i\omega t}}{P(i\omega)}
        \]
        Evaluamos $P(i\omega)$:
        \begin{align*}
            P(i\omega) &= (i\omega)^4 - 4(i\omega)^3 + 12(i\omega)^2 + 4(i\omega) - 13 \\
            &= \omega^4 + 4i\omega^3 - 12\omega^2 + 4i\omega - 13
        \end{align*}
        Agrupamos la parte real ($\alpha$) y la parte imaginaria ($\beta$):
        \[
            P(i\omega) = \underbrace{(\omega^4 - 12\omega^2 - 13)}_{\alpha} + i \underbrace{(4\omega^3 + 4\omega)}_{\beta} = \alpha + i\beta
        \]
        Sustituimos en nuestra solución compleja y multiplicamos por el conjugado para separar la parte real:
        \begin{align*}
            U_p &= \frac{\cos(\omega t) + i\sin(\omega t)}{\alpha + i\beta} \cdot \frac{\alpha - i\beta}{\alpha - i\beta} \\
            &= \frac{(\alpha \cos(\omega t) + \beta \sin(\omega t)) + i(\alpha \sin(\omega t) - \beta \cos(\omega t))}{\alpha^2 + \beta^2}
        \end{align*}
        La solución particular real $u_{p3}$ es la parte real de $U_p$:
        \[
            u_{p3} = \text{Re}(U_p) = \frac{\alpha \cos(\omega t) + \beta \sin(\omega t)}{\alpha^2 + \beta^2}
        \]

        \textbf{3. Solución general}

        Combinando todos los resultados, la solución general de la ecuación diferencial es:
        \[
            u(t) = c_1 e^t + c_2 e^{-t} + e^{2t}(c_3 \cos(3t) + c_4 \sin(3t)) - \frac{1}{13} + \frac{1}{27}e^{2t} + \frac{\alpha \cos(\omega t) + \beta \sin(\omega t)}{\alpha^2 + \beta^2}
        \]
        donde:
        \[
            \alpha = \omega^4 - 12\omega^2 - 13 \quad \text{y} \quad \beta = 4\omega^3 + 4\omega
        \]
        \item[(c)] 
        Resolvamos la ecuación diferencial:
        \[
            u^{(5)} + 3u^{(4)} + 3u''' + u'' = 1 + e^{2t}
        \]

        \textbf{1. Solución de la ecuación homogénea ($u_h$)}

        El polinomio característico asociado a la ecuación homogénea es:
        \[
            P(z) = z^5 + 3z^4 + 3z^3 + z^2
        \]
        Factorizamos extrayendo $z^2$ como factor común:
        \[
            P(z) = z^2 (z^3 + 3z^2 + 3z + 1)
        \]
        Reconocemos que el término entre paréntesis es el desarrollo de un cubo perfecto, $(z+1)^3$. Por lo tanto, el polinomio factorizado es:
        \[
            P(z) = z^2 (z+1)^3
        \]
        Las raíces del polinomio característico son:
        \begin{itemize}
            \item $z = 0$, raíz real doble (multiplicidad 2).
            \item $z = -1$, raíz real triple (multiplicidad 3).
        \end{itemize}
        Las funciones linealmente independientes asociadas a estas raíces forman la base del espacio nulo $\Sigma_0$:
        \[
            \mathcal{B} = \{ 1, t, e^{-t}, t e^{-t}, t^2 e^{-t} \}
        \]
        La solución general de la ecuación homogénea es:
        \[
            u_h(t) = c_1 + c_2 t + e^{-t}(c_3 + c_4 t + c_5 t^2)
        \]

        \textbf{2. Solución particular ($u_p$)}

        Aplicamos el principio de superposición para los términos del lado derecho $f(t) = 1 + e^{2t}$.

        \textit{Para el término constante $1$:}
        Normalmente propondríamos una constante $A$. Sin embargo, como $z = 0$ es raíz del polinomio característico con multiplicidad $s = 2$, la constante $1$ y la función $t$ ya forman parte de la solución homogénea. Por la regla de resonancia, debemos multiplicar nuestra propuesta por $t^2$:
        \[
            u_{p1} = A t^2
        \]
        Calculamos sus derivadas para sustituir en la ecuación:
        \[
            u_{p1}' = 2At, \quad u_{p1}'' = 2A, \quad u_{p1}''' = u_{p1}^{(4)} = u_{p1}^{(5)} = 0
        \]
        Sustituyendo en $P(D)u_{p1} = 1$:
        \[
            0 + 3(0) + 3(0) + 2A = 1 \implies 2A = 1 \implies A = \frac{1}{2}
        \]
        Por lo tanto, $u_{p1} = \frac{1}{2}t^2$.

        \textit{Para el término exponencial $e^{2t}$:}
        Como $2$ no es raíz del polinomio característico, evaluamos $P(2)$:
        \[
            P(2) = 2^2 (2+1)^3 = 4 \cdot 3^3 = 4 \cdot 27 = 108
        \]
        Proponemos $u_{p2} = B e^{2t}$ y usamos la propiedad $P(D)e^{\alpha t} = P(\alpha)e^{\alpha t}$:
        \[
            P(D)(B e^{2t}) = e^{2t} \implies B \cdot P(2) e^{2t} = e^{2t} \implies 108 B = 1 \implies B = \frac{1}{108}
        \]
        Por lo tanto, $u_{p2} = \frac{1}{108}e^{2t}$.

        La solución particular total es:
        \[
            u_p(t) = u_{p1} + u_{p2} = \frac{1}{2}t^2 + \frac{1}{108}e^{2t}
        \]

        \textbf{3. Solución general}

        La solución general de la ecuación diferencial es la suma de la solución homogénea y la solución particular ($u = u_h + u_p$):
        \[
            u(t) = c_1 + c_2 t + e^{-t}(c_3 + c_4 t + c_5 t^2) + \frac{1}{2}t^2 + \frac{1}{108}e^{2t}
        \]
        donde $c_1, c_2, c_3, c_4, c_5 \in \mathbb{R}$ son constantes arbitrarias.
    \end{enumerate}
}

% Ejercicio 32
\ejercicio{enunciado}{no}{
    Dado un entero $N \ge 2$ y $a_0, \dots, a_{N-1} \in \mathbb{C}$, considere la ecuación diferencial de orden $N$
    $$u^{(N)} = a_{N-1}u^{(N-1)} + \dots + a_1 u' + a_0 u,$$
    y suponga que $e^{\lambda_j t}$ es una solución para todo $j \in \{1, \dots, N\}$, con $\lambda_j \in \mathbb{C}$ y $\lambda_j \neq \lambda_k$ si $j \neq k$.
    \begin{enumerate}
        \item[(a)] Escriba un sistema de primer orden $U' = AU$, con $A \in \mathcal{M}_N(\mathbb{C})$, equivalente a la ecuación diferencial dada.
        \item[(b)] Determine el polinomio característico, los valores propios y los vectores propios de $A$.
        \item[(c)] Determine una matriz de soluciones de $U' = AU$.
        \item[(d)] Utilice la teoría de las Secciones 1.5 y 1.12 para demostrar que la matriz de Vandermonde
        $$\begin{pmatrix} 1 & \dots & 1 \\ \lambda_1 & \dots & \lambda_N \\ \vdots & \dots & \vdots \\ \lambda_1^{N-1} & \dots & \lambda_N^{N-1} \end{pmatrix}$$
        es invertible.
    \end{enumerate}
}{

}

% Ejercicio 33
\ejercicio{enunciado}{no}{
    Considere todas las soluciones de las ecuaciones diferenciales de este capítulo que no están definidas en todo $\mathbb{R}$. Analice el comportamiento de tales soluciones en la frontera de su dominio de definición.
}{

}

% Ejercicio 34
\ejercicio{enunciado}{no}{
    Suponga que
    $$A(t) = (a_{ij}(t))_{1 \le i, j \le N} \in \mathcal{M}_N(C([\alpha, \beta]; \mathbb{K})),$$
    con
    $$\text{tr } A(t) = 0, \quad \text{ para todo } t \in [\alpha, \beta].$$
    Para cada $t \in [\alpha, \beta]$, defina
    $$\Phi(t) : \mathbb{K}^N \to \mathbb{K}^N, \quad x \mapsto h(t; \alpha, x),$$
    donde $h(t; \alpha, x)$ denota la solución única de
    $$\begin{cases} h' = A(t)h, \\ h(\alpha) = x. \end{cases}$$
    Demuestre que, para cada $t \in [\alpha, \beta]$, $\Phi_t := \Phi(t)$ es un isomorfismo lineal de $\mathbb{K}^N$ que preserva el volumen, en el sentido de que, para cada conjunto medible $A \subset \mathbb{K}^N$,
    $$\text{vol}(A) := \int_A 1 \, dx = \int_{\Phi_t(A)} 1 \, dy =: \text{vol}(\Phi_t(A)).$$
}{

}

% Ejercicio 35
\ejercicio{enunciado}{no}{

}{

}

% Ejercicio 36
\ejercicio{enunciado}{no}{

}{

}

% Ejercicio 37
\ejercicio{enunciado}{no}{

}{

}

% Ejercicio 38
\ejercicio{enunciado}{no}{

}{

}


